% !TEX program = xelatex
\documentclass[a4paper]{article}
\usepackage[UTF8]{ctex}
\setCJKmainfont{SimSun}[BoldFont=SimHei,ItalicFont=KaiTi]
\setCJKsansfont{SimHei}
\usepackage[top=4mm,bottom=4mm,left=4mm,right=4mm]{geometry}
\usepackage{amsmath,amssymb,bm}
\usepackage{multicol}
\usepackage{graphicx}
\usepackage{tikz}
\usetikzlibrary{arrows.meta}
\usepackage[dvipsnames]{xcolor}
\usepackage{enumitem}
\usepackage{array,tabularx}
\pagestyle{empty}
\setlength{\parindent}{0pt}
\setlength{\parskip}{0.2pt}
\setlength{\columnsep}{2.8mm}
\setlength{\columnseprule}{0.15pt}
\renewcommand{\baselinestretch}{0.80}
\setlist{nosep,leftmargin=7pt,labelsep=1.5pt,itemsep=0pt,parsep=0pt,topsep=0pt}
\definecolor{h1}{HTML}{0D47A1}
\definecolor{h2}{HTML}{1565C0}
\definecolor{h3}{HTML}{E65100}
\definecolor{warnc}{HTML}{B71C1C}
\definecolor{tip}{HTML}{1B5E20}
\definecolor{bg}{HTML}{EEF2FF}
\definecolor{bgw}{HTML}{FFF3E0}
\newcommand{\chap}[1]{\vspace{1pt}{\fontsize{7.0}{7.3}\selectfont\bfseries\color{white}\colorbox{h1}{\strut\;#1\;}}\par\vspace{0.4pt}}
\newcommand{\sect}[1]{\vspace{0.5pt}{\fontsize{5.9}{6.2}\selectfont\bfseries\color{h2}$\blacktriangleright$ #1}\par}
\newcommand{\exm}[1]{{\fontsize{5.45}{5.8}\selectfont\bfseries\color{tip}$\diamond$ #1}\par}
\newcommand{\key}[1]{{\bfseries\color{h3}#1}}
\newcommand{\warn}[1]{{\bfseries\color{warnc}#1}}
\newcommand{\fml}[1]{\colorbox{bg}{$\displaystyle #1$}}
\newcommand{\wfml}[1]{\colorbox{bgw}{$\displaystyle #1$}}
\newcommand{\miniimg}[2]{\begin{center}\includegraphics[width=#1\columnwidth]{#2}\end{center}\vspace{-3pt}}
\newcommand{\schemabasics}{\begin{center}\begin{tikzpicture}[>=Latex,scale=0.58,every node/.style={font=\tiny}]
\draw[thick] (0,1.5)--(0,0)--(0.55,0)--(0.55,1.15);
\draw[thick] (1.2,1.35)--(1.2,0)--(0.65,0)--(0.65,0.95);
\draw[blue] (0.02,1.05)--(0.53,1.05) (0.67,0.82)--(1.18,0.82);
\node at (0.6,-0.25){U管/差压};
\draw[thick] (1.85,0.05)--(2.75,0.25)--(2.75,1.35)--(1.85,1.15)--cycle;
\draw[->,blue] (1.65,0.75)--(2.12,0.75);
\node at (2.35,1.55){闸门};
\draw[thick,->] (3.25,0.25)--(4.0,0.25)--(4.45,0.70);
\draw[thick] (3.25,0.45)--(3.95,0.45)--(4.30,0.80);
\draw[->,red] (3.10,0.35)--(3.45,0.35);
\draw[->,red] (4.25,0.65)--(4.58,0.98);
\node at (3.85,-0.25){弯管动量};
\draw[thick] (5.05,0.1)--(6.2,0.1);
\draw[thick,->] (5.25,0.25)--(6.05,0.25);
\node at (5.65,0.55){平板边界层};
\end{tikzpicture}\end{center}\vspace{-3pt}}
\newcommand{\schemepotential}{\begin{center}\begin{tikzpicture}[>=Latex,scale=0.58,every node/.style={font=\tiny}]
\draw[->] (-0.2,0)--(1.8,0) node[right]{$x$};
\draw[->] (0,-0.8)--(0,0.8) node[above]{$y$};
\foreach \y in {-0.45,0,0.45}{\draw[->,blue] (-0.05,\y)--(0.75,\y);}
\draw[thick] (1.05,0) circle (0.35);
\node at (1.05,-0.65){圆柱};
\draw[->,red] (2.35,0)--(3.1,0); \node at (2.25,0){源};
\draw[<- ,red] (3.8,0)--(4.55,0); \node at (4.65,0){汇};
\draw[thick] (2.75,-0.55).. controls (3.2,-0.9) and (3.75,-0.9)..(4.15,-0.55);
\draw[thick] (2.75,0.55).. controls (3.2,0.9) and (3.75,0.9)..(4.15,0.55);
\node at (3.45,-1.05){源汇/Rankine体};
\draw[thick] (5.2,-0.8)--(5.2,0.8);
\draw[fill=black] (5.75,0) circle (0.04); \node[right] at (5.78,0){涡};
\draw[fill=black] (4.65,0) circle (0.04); \node[left] at (4.62,0){镜像};
\node at (5.2,-1.05){壁面镜像};
\end{tikzpicture}\end{center}\vspace{-3pt}}
\newcommand{\schemelaval}{\begin{center}\begin{tikzpicture}[>=Latex,scale=0.58,every node/.style={font=\tiny}]
\draw[thick] (0,0.7).. controls (1.0,0.55) and (1.2,0.25)..(2.0,0.25).. controls (2.8,0.25) and (3.0,0.65)..(4.2,0.8);
\draw[thick] (0,-0.7).. controls (1.0,-0.55) and (1.2,-0.25)..(2.0,-0.25).. controls (2.8,-0.25) and (3.0,-0.65)..(4.2,-0.8);
\draw[dashed] (2.0,-0.55)--(2.0,0.55); \node at (2.0,-0.78){喉$A^*$};
\draw[->,blue] (-0.25,0)--(0.55,0); \node at (0,1.02){$p_0,T_0$};
\node at (4.15,1.02){$p_e,p_b$}; \draw[->,blue] (3.35,0)--(4.25,0);
\draw[red,thick] (3.0,-0.55)--(3.0,0.55); \node[red] at (3.22,0.72){正激波?};
\node at (2.2,1.12){Laval喷管};
\end{tikzpicture}\end{center}\vspace{-3pt}}
\newcommand{\schemeshock}{\begin{center}\begin{tikzpicture}[>=Latex,scale=0.58,every node/.style={font=\tiny}]
\draw[->,blue] (-0.2,0.45)--(1.0,0.45); \node at (0.4,0.7){$M_1>1$};
\draw[thick] (1.0,0)--(2.6,0);
\draw[thick] (1.0,0)--(2.6,0.55); \node at (2.25,0.78){楔角$\theta$};
\draw[red,thick] (1.0,0)--(1.75,1.05); \node[red] at (1.42,1.08){$\beta$};
\draw[->,blue] (1.9,0.36)--(2.75,0.65); \node at (2.65,0.35){$M_2$};
\draw[thick] (3.4,0)--(5.0,0); \draw[thick] (3.4,0)--(5.0,-0.55);
\foreach \a in {0.0,0.18,0.36,0.54}{\draw[red] (3.4,0)--(4.25,-\a);}
\node at (4.35,-0.85){膨胀扇$\delta$};
\end{tikzpicture}\end{center}\vspace{-3pt}}
\newcommand{\schemeviscous}{\begin{center}\begin{tikzpicture}[>=Latex,scale=0.58,every node/.style={font=\tiny}]
\draw[blue] (0,0)--(0,1.0); \draw[blue] (0.6,0)--(0.6,1.0);
\draw[thick] (0,1.0)--(0.6,1.0).. controls (1.2,1.0) and (1.2,0.3)..(1.9,0.3)--(3.3,0.3);
\draw[thick] (0,0)--(0.6,0).. controls (1.2,0) and (1.2,0.2)..(1.9,0.1)--(3.3,0.1);
\draw[->] (0.3,0.5)--(1.0,0.5); \node at (0.3,1.22){水面};
\node[circle,draw,inner sep=1pt] at (2.25,0.2){P};
\draw[<->] (2.25,0.2)--(2.25,-0.55); \node[right] at (2.25,-0.18){$h$};
\node at (1.65,-0.85){吸水泵/汽蚀};
\draw[thick] (4.05,0)--(5.75,0); \draw[thick] (4.05,0.65)--(5.75,0.65);
\draw[->,red] (4.1,0.65)--(5.2,0.65); \node at (5.35,0.83){$U$};
\draw[blue,thick] (4.1,0).. controls (4.8,0.25) and (5.2,0.48)..(5.75,0.65);
\node at (4.95,-0.25){Couette/Poiseuille};
\end{tikzpicture}\end{center}\vspace{-3pt}}
\newcommand{\schemebl}{\begin{center}\begin{tikzpicture}[>=Latex,scale=0.58,every node/.style={font=\tiny}]
\draw[thick] (0,0)--(3.4,0); \draw[->,blue] (-0.2,0.55)--(0.7,0.55); \node at (0.15,0.82){$U$};
\draw[red,thick] (0,0).. controls (1.1,0.25) and (2.4,0.55)..(3.4,0.85);
\draw[dashed] (2.6,0)--(2.6,0.68); \node[right] at (2.62,0.36){$\delta(x)$};
\node at (1.65,-0.25){平板边界层};
\draw[thick] (4.1,0) circle (0.65);
\draw[red,thick] (3.45,0.18).. controls (4.0,0.7) and (4.6,0.7)..(4.95,0.18);
\draw[red,thick] (3.45,-0.18).. controls (4.0,-0.7) and (4.6,-0.7)..(4.95,-0.18);
\node at (4.1,-0.95){圆柱分离尾流};
\end{tikzpicture}\end{center}\vspace{-3pt}}
\begin{document}
\fontsize{5.55}{6.15}\selectfont

% ==================== PAGE 1 ====================
\begin{multicols}{3}
\chap{P1 总判题地图：先选模型}
\schemabasics
\sect{考试资料策略}
可带：指定教材+计算器+3张A4双面。教材查长定义，本表放\key{题型识别、公式、答案骨架、易错点}。目标不是背诵，而是看到题能定位模型。

\sect{题干关键词定位}
\renewcommand{\arraystretch}{0.76}\begin{tabular}{|p{0.30\linewidth}|p{0.61\linewidth}|}
\hline 静止/水深/闸门/U管 & 静水压：$p=p_0+\rho gh$；平面/曲面总压力\\ \hline
喷管/Ma/背压/激波 & P3/P4：等熵、临界、激波、膨胀波\\ \hline
势函数/流函数/圆柱 & P2：$\varphi,\psi$叠加+伯努利\\ \hline
黏度/Re/管流/泵 & P5：N-S简化、H-P、管路损失\\ \hline
平板/边界层/阻力 & P6：$C_f,\delta,D$\\ \hline
弯管/喷流/受力 & 控制体动量：$\sum F=\dot m(V_2-V_1)$\\ \hline
\end{tabular}\renewcommand{\arraystretch}{1}

\sect{常数与基本量}
$g=9.81$；水$\rho=1000$；空气$\gamma=1.4,R=287$；$1\rm atm=101325Pa=10.33mH_2O$。\\
\fml{\nu=\mu/\rho,\quad Re=VL/\nu=\rho VL/\mu,\quad Ma=V/a,\quad a=\sqrt{\gamma RT}}

\sect{连续性}
不可压：$\nabla\cdot\bm V=0$；一维定常：$A_1V_1=A_2V_2=Q$。非均匀速度必须积分：$Q=\int_A v\,dA$。圆管抛物线层流：$\bar v=v_{\max}/2$。

\sect{控制体动量模板}
\fml{\sum F_x=\dot m(V_{2x}-V_{1x}),\quad \sum F_y=\dot m(V_{2y}-V_{1y})}, $\dot m=\rho Q$。\\
步骤：画CV；标速度方向；列压力/重力/支反力；动量写\key{流出-流入}；物体受力与流体受力反向。\\
\exm{弯管}
$x:p_1A_1-p_2A_2\cos\alpha-F_{Rx}=\rho Q(V_2\cos\alpha-V_1)$；\\
$y:-p_2A_2\sin\alpha+F_{Ry}=\rho QV_2\sin\alpha$。
\exm{射流冲板}
垂直板：$F=\rho QV$。斜板/分流：分$x,y$列动量，未给损失时各支速度近似$V$。

\sect{伯努利/机械能}
理想同流线：\fml{p/\rho g+V^2/(2g)+z=C}。\\
含损失/泵/水轮机：\\
\fml{\frac{p_1}{\rho g}+\frac{\alpha_1V_1^2}{2g}+z_1+H_p=\frac{p_2}{\rho g}+\frac{\alpha_2V_2^2}{2g}+z_2+h_f+\sum h_\zeta+H_T}
\warn{泵加左，水轮机输出加右；开口液面表压为0。}

\sect{静水压旧题急救}
平面：$F=\rho gh_CA$，$y_D=y_C+J_C/(y_CA)$。\\
曲面：$F_x=\rho gh_{Cx}A_x$；$F_y=\rho gV_p$，$V_p$为压力体体积；圆柱面合力过圆心。\\
U管爬楼梯：向下$+\rho gh$，向上$-\rho gh$，同一静止连通液体同高等压。
\sect{期末常见题一眼归类}
\renewcommand{\arraystretch}{0.70}\begin{tabular}{|p{0.23\linewidth}|p{0.69\linewidth}|}
\hline 压力中心 & 先算形心压强总力，再用惯性矩修正作用点。\\ \hline
曲面受压 & 水平分量投影面；竖直分量为压力体重量。\\ \hline
流线/迹线 & 定常二者重合；非定常迹线要解质点方程。\\ \hline
量纲分析 & 列变量，选重复变量，写$\pi$群，最后说明物理意义。\\ \hline
管路/泵 & 永远从大水面、压力表、出口三类点中选两点列能量。\\ \hline
\end{tabular}\renewcommand{\arraystretch}{1}
\sect{无脑列方程顺序}
1 定研究对象：质点/控制体/流线/截面。2 写主方程：连续、动量、能量、N--S、状态方程。3 写边界条件。4 单位统一。5 结果检查：压强是否表压、速度是否平均、面积是否截面。

\sect{概念判断}
亚声速椭圆型、全局影响；超声速双曲型、马赫锥内传播。激波绝热非等熵，$T_0$不变、$p_0$下降。膨胀波近似等熵。逆压梯度易分离。
\sect{速度场题}
迹线：$dx/dt=u,dy/dt=v,dz/dt=w$，给初始点积分。流线：定常时$dx/u=dy/v=dz/w$。加速度：$\bm a=\partial\bm V/\partial t+u\bm V_x+v\bm V_y+w\bm V_z$。
\sect{旋度/应变/应力}
$\nabla\cdot\bm V=u_x+v_y+w_z$。二维旋度$\Omega_z=v_x-u_y$，角速度常取$\omega_z=\Omega_z/2$。牛顿流体：$\tau_{xy}=\mu(u_y+v_x)$；不可压正应力粘性项$2\mu u_x,2\mu v_y,2\mu w_z$。
\sect{静水压补充模板}
倾斜平板：$h_C=\bar y\sin\theta$，$F=\rho gh_CA$，压力中心沿板坐标$y_D=y_C+I_C/(y_CA)$。相对平衡：水平加速度$a$时自由面$\tan\alpha=a/g$；旋转液面$z=\omega^2r^2/(2g)+C$。
\sect{量纲分析模板}
设$f(x_1,\cdots,x_n)=0$，选$k$个重复变量覆盖基本量纲，构造$n-k$个无量纲数$\pi_i$。常用：$Re=\rho VL/\mu$，$Fr=V/\sqrt{gL}$，$Eu=\Delta p/(\rho V^2)$，$We=\rho V^2L/\sigma$，$Ma=V/a$。
\sect{控制体符号检查}
压力力永远指向控制体内；出口速度按流出方向取正。若题问管壁受力，先求“壁对流体”，最后取反。开口大水面$V\approx0$，表压$p=0$。
\sect{应力张量题}
给应力张量$P$和平面法向$\bm n$：应力矢量$\bm p_n=P\bm n$。法向分量$p_N=\bm p_n\cdot\bm n$；切向矢量$\bm p_T=\bm p_n-p_N\bm n$；夹角$\cos\alpha=p_N/|\bm p_n|$。圆柱面法向用几何外法线。
\sect{平面受压常用惯性矩}
矩形宽$b$高$h$：$I_C=bh^3/12$；圆面积：$I_C=\pi R^4/4$。压力中心永远比形心更深。若要求铰链力，先算液压力合力和作用点，再对铰链取矩。
\sect{小孔/水箱出流}
大水箱小孔：$V=C_v\sqrt{2gH}$，$Q=C_dA\sqrt{2gH}$。水位下降：$-A_t\,dH/dt=C_dA_o\sqrt{2gH}$，积分得时间。若是水银压差计，先换算实际压差。
\sect{U管振荡}
等截面U管总液柱长$L$，位移$z$时回复力$2\rho gAz$，质量$\rho AL$，故$\ddot z+(2g/L)z=0$，$T=2\pi\sqrt{L/(2g)}$。有非定常伯努利也可推同式。
% ==================== SECTION 2 ====================
\chap{P2 势流叠加：理想不可压无旋}
\schemepotential
\sect{基本条件}
无旋：$\nabla\times\bm V=0\Rightarrow \bm V=\nabla\varphi$。二维不可压可设流函数$\psi$。\\
\fml{\nabla^2\varphi=0,\quad \nabla^2\psi=0}
直角：$u=\varphi_x=\psi_y,\;v=\varphi_y=-\psi_x$。\\
极坐标：\fml{v_r=\varphi_r=\frac1r\psi_\theta,\quad v_\theta=\frac1r\varphi_\theta=-\psi_r}

\sect{基本流表}
\renewcommand{\arraystretch}{0.76}\begin{tabular}{|p{0.18\linewidth}|p{0.35\linewidth}|p{0.35\linewidth}|}
\hline 流动 & $\varphi$ & $\psi$\\ \hline
均匀流 & $Ur\cos\theta$ & $Ur\sin\theta$\\ \hline
源 & $\frac{Q}{2\pi}\ln r$ & $\frac{Q}{2\pi}\theta$\\ \hline
汇 & $-\frac{Q}{2\pi}\ln r$ & $-\frac{Q}{2\pi}\theta$\\ \hline
点涡 & $\frac{\Gamma}{2\pi}\theta$ & $-\frac{\Gamma}{2\pi}\ln r$\\ \hline
偶极子 & $\frac{M\cos\theta}{r}$ & $-\frac{M\sin\theta}{r}$\\ \hline
\end{tabular}\renewcommand{\arraystretch}{1}
\warn{叠加法：总$\varphi,\psi$直接相加，速度也相加。}
\sect{看到题先搭哪几个基元}
\renewcommand{\arraystretch}{0.70}\begin{tabular}{|p{0.28\linewidth}|p{0.62\linewidth}|}
\hline 均匀来流绕物体 & 均匀流+偶极子；有升力再加点涡。\\ \hline
半无限体/Rankine体 & 均匀流+点源；分界流线取过驻点的$\psi$。\\ \hline
壁面附近源/涡 & 用镜像让壁面成为流线：源同号、涡异号。\\ \hline
边角绕流 & 用$\varphi=Ar^n\cos n\theta$，$n=\pi/\alpha$。\\ \hline
\end{tabular}\renewcommand{\arraystretch}{1}

\sect{给速度求$\varphi,\psi$}
有$\varphi$先验：$v_x-u_y=0$；有$\psi$先验：$u_x+v_y=0$。\\
$\varphi_x=u\Rightarrow \varphi=\int u dx+f(y)$，再用$\varphi_y=v$定$f$。\\
$\psi_y=u\Rightarrow \psi=\int u dy+g(x)$，再用$-\psi_x=v$定$g$。

\sect{圆柱绕流}
\fml{\varphi=U(r+a^2/r)\cos\theta,\quad \psi=U(r-a^2/r)\sin\theta}\\
$v_r=U(1-a^2/r^2)\cos\theta$，$v_\theta=-U(1+a^2/r^2)\sin\theta$。\\
壁面$r=a$：$v_r=0,\;v_\theta=-2U\sin\theta$；驻点$\theta=0,\pi$；最大速$2U$在$90^\circ,270^\circ$。\\
压强：\fml{p=p_\infty+\frac12\rho(U^2-V^2)}。壁面$p=p_\infty+\frac12\rho U^2(1-4\sin^2\theta)$。
\sect{圆柱绕流答题模板}
写“均匀流+偶极子叠加，$r=a$为$\psi=0$流线”。若求压力分布，先求$V^2=v_r^2+v_\theta^2$；若求合力，压力积分：$dF_x=-p\cos\theta a\,d\theta$，$dF_y=-p\sin\theta a\,d\theta$。无环量阻力升力均为0；有环量只剩升力$\rho U\Gamma$。

\sect{有环量/升力}
圆柱+点涡：$v_\theta=-2U\sin\theta+\Gamma/(2\pi a)$。升力：\fml{L'=\rho U\Gamma}。理想无粘绕流阻力$D=0$。

\sect{镜像法/边角}
固壁为流线，要求法向速度为0。源/汇靠直壁：同号镜像；点涡靠直壁：反号镜像。半平面源汇常见$2\pi\to\pi$。\\
边角：若$\varphi=Ar^n\cos n\theta,\psi=Ar^n\sin n\theta$，壁面$\psi=0$，夹角$\alpha$满足$n=\pi/\alpha$。$V\sim r^{n-1}$。
\sect{源汇/驻点常用算式}
源+均匀流：$v_r=U\cos\theta+Q/(2\pi r)$，$v_\theta=-U\sin\theta$。驻点在$\theta=\pi$，$r=Q/(2\pi U)$。过驻点分界流线$\psi=Q/2$。源汇对：两个源汇距离给出时，写总$\psi$后令$\psi=C$画物体轮廓；速度由偏导，不要凭图猜。

\sect{常考势流原题骨架}
\exm{均匀流+点源半体}
$\psi=Ur\sin\theta+Q\theta/(2\pi)$。驻点$\theta=\pi,r=Q/(2\pi U)$；分界线$\psi=Q/2$。远处半体宽度$b=Q/U$。
\exm{均匀流+源汇对}
源在$(-a,0)$、汇在$(a,0)$：总$\psi=Uy+\frac{Q}{2\pi}(\theta_1-\theta_2)$。驻点由$u=v=0$求；封闭流线为Rankine卵形。
\exm{壁面吸水/缝隙}
半平面源汇常从$2\pi$改成$\pi$；先写$\psi=Uy-Q\theta/\pi$，令壁面$\psi=0$，驻点由速度为0定位置。
\exm{圆柱有环量停滞点}
壁面$v_\theta=-2U\sin\theta+\Gamma/(2\pi a)$。停滞点满足$\sin\theta=\Gamma/(4\pi Ua)$；若绝对值$>1$，停滞点离开圆柱表面。
\sect{伯努利用法}
势流定常无粘：$p+\frac12\rho V^2=p_\infty+\frac12\rho U^2$。求压力最低点先找$V$最大点；求空化/吸力，比较$p$与汽化压或外界压。
\sect{给速度场求势/流函数例式}
若$u=3x^2-3y^2,\;v=-6xy$：$u_y=-6y,v_x=-6y$，无旋；$\varphi=\int u dx=x^3-3xy^2+C$。$\psi_y=u\Rightarrow \psi=3x^2y-y^3+g(x)$，由$-\psi_x=v$得$g'=0$。
\sect{源/涡压力题}
点涡$V_\theta=\Gamma/(2\pi r)$，伯努利给$p(r)=C-\frac12\rho\Gamma^2/(4\pi^2r^2)$，越靠中心压强越低。源汇有径向速度，压力同样按$V^2$比较。
\sect{圆柱绕流带答案}
无环量：$C_p=(p-p_\infty)/(0.5\rho U^2)=1-4\sin^2\theta$。有环量：把$v_\theta$代入$C_p=1-(v_\theta/U)^2$。升力方向按环量和来流右手关系判。
\sect{保角/边角题}
角域$\alpha$内势流解常取$n=\pi/\alpha$。若$\alpha<\pi$，$n>1$，角点速度趋0；若$\alpha>\pi$，$n<1$，角点速度趋无穷，物理上需粘性修正。

\sect{答案句}
“该流动不可压无旋，故$\varphi,\psi$满足Laplace方程。由叠加法写总流函数，速度由极坐标关系求得，固壁为$\psi=C$。压强由伯努利求。”
% ==================== SECTION 3 ====================
\chap{P3 等熵流 + Laval 喷管}
\schemelaval
\sect{基本公式}
\fml{p=\rho RT,\quad a=\sqrt{\gamma RT},\quad Ma=V/a}
\fml{\frac{T_0}{T}=1+\frac{\gamma-1}{2}M^2}
\fml{\frac{p_0}{p}=\left(1+\frac{\gamma-1}{2}M^2\right)^{\gamma/(\gamma-1)}}
\fml{\frac{\rho_0}{\rho}=\left(1+\frac{\gamma-1}{2}M^2\right)^{1/(\gamma-1)}}\\
空气：$T_0/T=1+0.2M^2$，$p_0/p=(1+0.2M^2)^{3.5}$，$\rho_0/\rho=(1+0.2M^2)^{2.5}$。

\sect{临界/阻塞}
$M=1,A=A^*$。\\
\fml{T^*/T_0=0.833,\quad p^*/p_0=0.528,\quad \rho^*/\rho_0=0.634}\\
\warn{0.528只表示临界静压比，不是所有出口/背压比。}
\sect{总量、静量、背压别混}
$p_0,T_0$是把当地气体等熵减速到$M=0$得到的量；等熵段不变，过激波$p_0$降、$T_0$不变。$p_e$是喷管出口截面内静压，$p_b$是外界背压。收缩管未阻塞时$p_e=p_b$；Laval非设计时出口外还会用激波/膨胀波调整。

\sect{面积--马赫数}
\fml{\frac{A}{A^*}=\frac1M\left[\frac{2}{\gamma+1}\left(1+\frac{\gamma-1}{2}M^2\right)\right]^{\frac{\gamma+1}{2(\gamma-1)}}}\\
$A/A^*>1$通常有亚声速/超声速两支。收缩管最多到$M=1$；Laval扩张段可超声速。

\sect{质量流量}
\fml{\dot m=\rho VA=ApM\sqrt{\gamma/(RT)}}\\
阻塞最大流量在$M=1$，空气：\fml{\dot m_{\max}=0.0404A^*p_0/\sqrt{T_0}}。
\sect{Ma作为中间变量}
知道某截面$M$，即可由等熵式求该截面$T,p,\rho,V$；但跨激波不能直接连，要先用激波关系换到波后，再以新的$p_{02}$继续等熵。

\sect{Laval 背压状态}
\renewcommand{\arraystretch}{0.76}\begin{tabular}{|p{0.30\linewidth}|p{0.61\linewidth}|}
\hline 背压高 & 全管亚声速，$p_e=p_b$，未阻塞\\ \hline
喉部$M=1$ & 阻塞，$\dot m$最大\\ \hline
扩张段正激波 & 波前超声速，波后亚声速，$p_0$降\\ \hline
设计状态 & 全管等熵，$p_e=p_b$\\ \hline
$p_e<p_b$ & 过膨胀，出口外压缩/斜激波\\ \hline
$p_e>p_b$ & 欠膨胀，出口外膨胀波\\ \hline
\end{tabular}\renewcommand{\arraystretch}{1}

\sect{收缩喷管题}
给$p_0,T_0,p_b,A_e$。若$p_b/p_0>0.528$：未阻塞，$p_e=p_b$，用$p_0/p_e$求$M_e$，再求$T_e,\rho_e,V_e,\dot m=\rho_eV_eA_e$。\\
若$p_b/p_0\le0.528$：阻塞，出口/喉部$M=1$，$\dot m=0.0404A_ep_0/\sqrt{T_0}$。
\sect{收缩管7.5标准步骤}
入口速度可忽略：入口给的$p,T$就是$p_0,T_0$。1 算$p_b/p_0$。2 与0.528比。3 未阻塞：令$p_e=p_b$求$M_e$，再$\rho_e=p_e/(RT_e)$，$\dot m=\rho_eM_e\sqrt{\gamma RT_e}A_e$。4 阻塞：$M_e=1$，$p_e=p^*$不是$p_b$，用最大流量式。

\sect{Laval面积/流量题}
给$A^*$与出口$p_e$或$T_e$，无激波。用等熵求$M_e$；若取超声速支，则喉部$M=1,A_t=A^*$。出口面积由$A_e/A^*$求；流量可用喉部最大流量或出口$\rho VA$。
\sect{Laval7.6标准步骤}
1 由出口$T_e$或$p_e$和$p_0,T_0$反求$M_e$。2 题说无激波且$M_e>1$，出口必须取超声速支。3 由$A_e/A^*(M_e)$求$A_e$。4 流量用$\dot m_{\max}$或出口$\rho_eV_eA_e$；两者应一致。

\sect{文丘里质量流量计}
给最大流量和最低喉压$p_t$。先由入口$M_1$求$p_0,T_0$，再用$p_0/p_t=(1+0.2M_t^2)^{3.5}$求$M_t$。若$M_t<1$，不可用0.528，继续算$T_t,\rho_t,V_t,A_t$。
\sect{文丘里7.7标准步骤}
最大流量$\dot m$和入口面积给出：先用$\dot m=\rho_1V_1A_1$求$V_1,M_1$，再求$p_0,T_0$。最低允许喉压$p_t$给出后求$M_t$。若$M_t=1$才临界；通常$M_t<1$，所以不能写$p_t/p_0=0.528$。
\sect{可压缩题判断树}
1 题给$p,T,V$：先算$M$，再算$p_0,T_0$。2 题给$p_0,T_0$和某静压：用$p_0/p$反求$M$。3 题给面积比：用$A/A^*$反求$M$，并按题意选亚声速/超声速支。4 涉及背压：先判断是否阻塞，再判断是否有激波。
\sect{几个常用数值}
空气$M=1$：$T/T_0=0.833,p/p_0=0.528,\rho/\rho_0=0.634$。$M=2$：$T_0/T=1.8,p_0/p\approx7.82$。$M=3$：$T_0/T=2.8,p_0/p\approx36.7$。
\sect{答题检查句}
“本段若无摩擦和激波，按一维绝热等熵处理，总温总压不变；若出现正激波，总温仍不变、总压下降，激波后应使用新的滞止压力继续计算。”
\sect{可压缩基础题}
弹性模量$E=-dp/(dV/V)$，密度相对变化$d\rho/\rho=dp/E$。声速$a=\sqrt{K/\rho}$或理想气$a=\sqrt{\gamma RT}$。飞机声响题：若$M=V/a>1$，马赫锥半角$\sin\mu=1/M$，听到时间/距离由几何三角形算。
\sect{喷管出口状态三句话}
收缩喷管：出口就是最小面积，阻塞时出口$M=1$。Laval喷管：喉部是最小面积，阻塞时喉部$M=1$，出口可亚声速或超声速。背压只直接等于出口静压的情况：未阻塞或完全匹配设计状态。
\sect{面积比求M查表替代}
考试无表时，写“由面积--马赫数关系迭代解得$M=\cdots$”。若$A/A^*$给定且题说超声速，必须取大于1那支；若收缩管，不能取超声速支。
\sect{质量流量检查}
同一截面三种写法应一致：$\dot m=\rho VA=ApM\sqrt{\gamma/(RT)}=A p_0\sqrt{\gamma/(RT_0)}M(1+0.2M^2)^{-3}$（空气）。阻塞时$M=1$代入得$0.0404A^*p_0/\sqrt{T_0}$。
% ==================== SECTION 4 ====================
\chap{P4 激波 + 膨胀波}
\schemeshock
\sect{总判断}
超声速遇压缩角/凹角/迎风面$\Rightarrow$激波；遇扩张角/凸角$\Rightarrow$膨胀波；偏转角太大$\Rightarrow$脱体激波。激波：$p,T,\rho\uparrow$，$V,M\downarrow$，$T_0$不变，$p_0$下降。

\sect{正激波}
\fml{M_2^2=\frac{1+\frac{\gamma-1}{2}M_1^2}{\gamma M_1^2-\frac{\gamma-1}{2}}}
空气：$M_2^2=(1+0.2M_1^2)/(1.4M_1^2-0.2)$。\\
\fml{\frac{p_2}{p_1}=1+\frac{2\gamma}{\gamma+1}(M_1^2-1)}
\fml{\frac{\rho_2}{\rho_1}=\frac{(\gamma+1)M_1^2}{(\gamma-1)M_1^2+2}}
\fml{\frac{T_2}{T_1}=(p_2/p_1)/(\rho_2/\rho_1)}\\
$p_{02}/p_{01}=(p_2/p_1)(p_{02}/p_2)/(p_{01}/p_1)$，用等熵式算$p_0/p$。

\sect{斜激波流程}
$M_{n1}=M_1\sin\beta$。本质：法向分量过正激波，切向分量不变。\\
已知$M_1,\theta$：查$\theta-\beta-M$表得$\beta$；算$M_{n1}$；用正激波公式得$p_2/p_1,T_2/T_1,M_{n2}$；$M_2=M_{n2}/\sin(\beta-\theta)$。工程常取弱激波解。
\sect{斜激波无表时}
\fml{\tan\theta=2\cot\beta\frac{M_1^2\sin^2\beta-1}{M_1^2(\gamma+\cos2\beta)+2}}\\
先数值/试算$\beta$。若偏转角超过最大值，无附体斜激波，写“脱体激波”，不能强行用弱解。

\sect{膨胀波}
\fml{\nu(M)=\sqrt{\frac{\gamma+1}{\gamma-1}}\tan^{-1}\sqrt{\frac{\gamma-1}{\gamma+1}(M^2-1)}-\tan^{-1}\sqrt{M^2-1}}\\
空气：$\nu(M)=\sqrt6\tan^{-1}\sqrt{(M^2-1)/6}-\tan^{-1}\sqrt{M^2-1}$。\\
膨胀角：\fml{\delta=\nu(M_2)-\nu(M_1)}。给$\delta$：$\nu(M_2)=\nu(M_1)+\delta$。
\sect{膨胀波后物理量}
膨胀波近似等熵：$p_0,T_0$不变。已知$M_1,M_2$后：
$p_2/p_1=\left[\frac{1+0.2M_1^2}{1+0.2M_2^2}\right]^{3.5}$，
$T_2/T_1=(1+0.2M_1^2)/(1+0.2M_2^2)$。

\sect{7.19 管外膨胀模板}
给$A_e/A_t,p_b,p_0$。1 由$A_e/A^*$取超声速支求$M_e$。2 等熵求$p_e=p_0/(1+0.2M_e^2)^{3.5}$。3 若$p_e>p_b$，出口外继续膨胀。4 由$p_0/p_b$求最终$M_b$。5 $\alpha=\nu(M_b)-\nu(M_e)$。

\sect{7.20 管内正激波模板}
给激波位置$A_s/A_t$。1 激波前由$A_s/A^*$超声速支求$M_1$。2 正激波求$M_2,p_2/p_1,p_{02}/p_{01}$。3 激波后以新总压$p_{02}$亚声速等熵到出口。
\sect{Laval内激波完整链}
$p_{01}=p_0$。波前：$A_s/A_t\to M_1\to p_1,T_1$。过波：$M_1\to M_2,p_2,T_2,p_{02}$。波后：用$A_e/A_s$或等效$A^*_2$求出口$M_e$，再用$p_{02}/p_e$求$p_e$，最后令$p_b=p_e$判断该激波位置是否匹配。

\sect{易错}
知道$M_2$不能随便算全部：激波后静量接激波前$p_1,T_1,\rho_1$。$p_{01}$是波前滞止压，$p_{02}$是波后滞止压，$p_{02}<p_{01}$。膨胀波不用正激波公式。
\sect{激波反射/相交题写法}
把每道波按顺序编号。每过一道斜激波：$M,p,T,\rho$用斜激波关系更新；每过一道膨胀波：用PM函数更新$M$，再等熵更新$p,T$。两股流相交后的接触面两侧静压相等，方向用偏转角闭合。
\sect{压缩角/膨胀角判断}
超声速来流沿壁面走：壁面向流体内折，流体被压缩$\Rightarrow$斜激波；壁面向外展开，流体膨胀$\Rightarrow$膨胀扇。亚声速不能套斜激波/PM函数。
\sect{Laval背压八状态速记}
背压从高到低：不流动$\to$全亚声速$\to$喉部临界$\to$扩张段内正激波$\to$出口正激波$\to$出口外斜激波(过膨胀)$\to$设计等熵$\to$出口外膨胀波(欠膨胀)。质量流量一旦阻塞后基本不再随背压增加。
\sect{写题保底句}
“激波前后满足质量、动量、能量守恒，但熵增，故不可用同一个$p_0$贯穿波前波后。”这句话通常能救概念分。
\sect{正激波常用极限}
$M_1\to1$时波很弱，$M_2\to1$、总压损失小。$M_1$很大时，空气$\rho_2/\rho_1\to6$。正激波后一定亚声速，斜激波后可亚声速也可超声速。
\sect{斜激波压力比}
只要算出$M_{n1}$，压力比直接用正激波式：$p_2/p_1=1+\frac{2\gamma}{\gamma+1}(M_{n1}^2-1)$。温度、密度同理。不要把$M_1$直接代正激波式。
\sect{膨胀波常见题}
给$M_1$和转角$\delta$：查/算$\nu(M_1)$，得$\nu(M_2)=\nu(M_1)+\delta$，反查$M_2$。压力用等熵比；若要求壁面角，壁面转过的角就是膨胀角。
\sect{管外膨胀7.19简答句}
若由面积比得到的等熵出口压$p_e$大于背压$p_b$，说明出口射流继续膨胀，属于管外膨胀；最终射流压力调至$p_b$，角度用PM函数差。
% ==================== SECTION 5 ====================
\chap{P5 粘性管流：N-S、H-P、Couette、管路}
\schemeviscous
\sect{N-S 简化}
不可压牛顿流体：$\rho D\bm V/Dt=-\nabla p+\mu\nabla^2\bm V+\rho\bm g$。定常$\partial_t=0$；单向流只保留$u$；充分发展$\partial u/\partial x=0$；平行流惯性项常为0。\\
\warn{大题先写假设，再积分两次，再代边界条件。}

\sect{H-P 圆管层流}
\fml{u(r)=\frac{1}{4\mu}\left(-\frac{dp}{dx}\right)(R^2-r^2)}
\fml{u_{\max}=2\bar u,\quad Q=\frac{\pi R^4}{8\mu}\left(-\frac{dp}{dx}\right)}
\fml{\Delta p=\frac{8\mu LQ}{\pi R^4},\quad \lambda=64/Re}
壁剪：$\tau_w=R\Delta p/(2L)=8\mu\bar u/d$。
\sect{广义压强写法}
有重力沿管时，用$p^*=p+\rho gz$。H--P可写成$u(r)=\frac{1}{4\mu}(-dp^*/dx)(R^2-r^2)$。斜管、竖管不要漏掉$\rho g$项。

\sect{Couette 流}
无压差：$u=Uy/h,\;\tau=\mu U/h$。\\
有压力梯度/重力分量：\fml{\mu d^2u/dy^2=dp/dx-\rho g_x}，积分两次。\\
倾斜平板若$u(0)=0,u(H)=-U_0,u(H/2)=0$，且$\mu u''+\rho g\sin\theta=0$，则$U_0=\rho g\sin\theta H^2/(4\mu)$。

\sect{双层 Couette}
\fml{\tau=\frac{U}{h_1/\mu_1+h_2/\mu_2}=\frac{U\mu_1\mu_2}{\mu_2h_1+\mu_1h_2}}
$u_1=\tau y/\mu_1$；$u_2=u_1(h_1)+\tau(y-h_1)/\mu_2$。

\sect{管路损失}
\fml{h_f=\lambda\frac{L}{d}\frac{V^2}{2g},\quad h_\zeta=\zeta\frac{V^2}{2g}}
粗糙度只通过$\varepsilon/d$影响$\lambda$；层流$\lambda=64/Re$与粗糙度无关。\\
\fml{\frac1{\sqrt\lambda}=-2\log_{10}\left(\frac{\varepsilon}{3.7d}+\frac{2.51}{Re\sqrt\lambda}\right)}\\
无Moody图时用Colebrook试算；完全粗糙区近似只看$\varepsilon/d$。\\
吸水泵：$h=(p_a-p_B)/\rho g-V^2/2g-(\lambda L/d+\sum\zeta)V^2/2g$。\\
水轮机：$H_T=z_1-z_2+(p_1-p_2)/\rho g+(V_1^2-V_2^2)/2g-h_f-\sum h_\zeta$，$P=\rho gQH_T$。
\sect{管路题统一流程}
1 由$Q$算各段$V=4Q/(\pi d^2)$。2 算$Re=Vd/\nu$判层/湍。3 层流$\lambda=64/Re$；湍流用Moody/题给提示。4 各段$h_f$相加，弯头阀门用$\sum h_\zeta$。5 能量方程求$p,H_p,H_T$。

\sect{常见题模板}
\exm{泵最大安装高度}
算$Q,V,Re,\varepsilon/d\to\lambda$；入口压力要求$p_B\ge p_v+\Delta p_{\text{safe}}$；代入求$h_{\max}$。
\exm{水轮机功率}
各段算$Re,\lambda,h_f$；能量方程含$H_T$；$P=\rho gQH_T$。
\exm{滑动轴承/油膜}
$du/dy\approx V/t$；$\tau=\mu V/t$；$F=\tau A$；$P=FV$。
\exm{题2吸水泵高度}
起点选水面，终点选泵入口中心。入口阀、弯头、沿程损失全在右侧；泵入口最低绝压$p_B=p_v+20{\rm kPa}$。粗糙度用于$\varepsilon/d$查$\lambda$。
\exm{题3水轮机输出}
从上游水面到下游压力表截面列能量；不同直径分段算$V,Re,\lambda,h_f$。压力表给表压，水面表压为0；输出功率$P=\rho gQH_T$。
\sect{平板Poiseuille}
两固定板距$h$，定常充分发展：$\mu u''=dp/dx$。边界$u(0)=u(h)=0$，得$u=\frac{1}{2\mu}(-dp/dx)y(h-y)$，$u_{\max}=\frac{h^2}{8\mu}(-dp/dx)$，$\bar u=2u_{\max}/3$。
\sect{Couette--Poiseuille合成}
上板$U$、下板0、同时有压差：$u=Uy/h+\frac{1}{2\mu}(-dp/dx)y(h-y)$。零剪切位置由$du/dy=0$；流量$Q=b\int_0^h udy$。
\sect{层流/湍流判定}
圆管$Re<2300$层流，$Re>4000$通常湍流。非圆管用水力直径$d_h=4A/P_w$。层流有解析速度分布；湍流多用经验阻力系数，不能套H--P抛物线。
\sect{压力降与剪力平衡}
圆管受力平衡：$\Delta p\,\pi R^2=\tau_w(2\pi RL)$，故$\tau_w=R\Delta p/(2L)$。这条可推很多壁面剪应力题。
\sect{管内层流竖直流}
沿$z$向上流，能驱动流动的是$-\frac{dp}{dz}-\rho g$。若速度分布已知，由H--P反推$dp/dz$；题问“沿轴向压降”要说明是否含重力项，常用$p^*=p+\rho gz$避免漏项。
\sect{第8章层流作业套路}
两平行板/两层流体：主方程都是$\mu u''=$常数。边界：固壁无滑移；两流体交界速度连续、剪应力连续。先解速度，再求剪应力$\tau=\mu u'$，最后代数值。
\sect{泵/汽蚀判断}
入口绝压必须大于汽化压加安全余量。若算得$p_B<p_v+\Delta p$，就汽蚀。吸水高度越大、流速越大、管越长、局部损失越大，越容易汽蚀。
\sect{局部损失系数}
入口、弯头、阀门、出口统一写$h_\zeta=\zeta V^2/(2g)$。若不同管径，局部损失用该局部元件所在管段速度，不能全用一个$V$。
% ==================== SECTION 6 ====================
\chap{P6 边界层、阻力、最后急救}
\schemebl
\sect{平板边界层}
$Re_x=Ux/\nu,\;Re_L=UL/\nu$。层流：\\
\fml{\delta=5x/\sqrt{Re_x},\quad \bar C_f=1.328/\sqrt{Re_L}}\\
湍流：\fml{\bar C_f=0.074/Re_L^{1/5}}；转捩：\fml{\bar C_f=0.074/Re_L^{1/5}-1742/Re_L}。\\
\fml{D=\bar C_f\frac12\rho U^2A,\quad P=DU}

\sect{位移/动量厚度}
\fml{\delta^*=\int_0^\delta(1-u/U)dy,\quad \theta=\int_0^\delta(u/U)(1-u/U)dy}
\fml{\tau_w/(\rho U^2)=d\theta/dx}
给$u/U=f(\eta),\eta=y/\delta$：把$dy=\delta d\eta$代入积分。
\sect{常见剖面直接用}
线性$u/U=\eta$：$\delta^*=\delta/2,\theta=\delta/6$。抛物$u/U=2\eta-\eta^2$：$\delta^*=\delta/3,\theta=2\delta/15$。若题给幂律$u/U=\eta^{1/n}$，先代$\eta$积分，不要展开成$y$。

\sect{书10.2 平板油流}
给$\mu,\rho,U,L,b$。先$\nu=\mu/\rho$，$Re_L=UL/\nu$。若层流：$\delta_{\max}=5L/\sqrt{Re_L}$，两面阻力$D=\bar C_f(0.5\rho U^2)(2Lb)$。\\
\warn{这是平板边界层，不是管流壁律。}
\sect{平板题统一流程}
1 算$Re_L$判层/湍/转捩。2 选$\bar C_f$和$\delta(L)$。3 阻力面积：单面$A=Lb$，双面$A=2Lb$。4 题问尾部边界层厚度用$x=L$；问总阻力用平均摩擦系数。

\sect{湍流近壁三层}
$u_\tau=\sqrt{\tau_w/\rho}$，$y^+=yu_\tau/\nu$，$u^+=\bar u/u_\tau$。线性底层$y^+<5,u^+=y^+$；缓冲层$5<y^+<30$；湍流核心$u^+=2.5\ln y^+ +5.5$。圆管平均：$\bar u/u_\tau=2.5\ln(Ru_\tau/\nu)+1.75$。
\sect{壁律题写法}
给管径、平均速度、$\rho,\mu$：先$\nu=\mu/\rho$。由$\bar u/u_\tau=2.5\ln(Ru_\tau/\nu)+1.75$试算$u_\tau$，再$\tau_w=\rho u_\tau^2$。线性底层最大厚度：$y=5\nu/u_\tau$；缓冲层到$y=30\nu/u_\tau$。

\sect{分离与阻力}
逆压梯度$dp/dx>0$使边界层减速增厚；$\partial u/\partial y|_w=0$为分离临界。摩擦阻力来自$\tau_w$；压差阻力来自分离尾流。
\sect{外绕流阻力判断}
流线型小迎角：摩擦阻力为主；钝体/分离明显：压差阻力为主。高$Re$物面外可近似无粘，用Euler/伯努利；靠壁必须用边界层。

\sect{最后急救题库}
\renewcommand{\arraystretch}{0.76}\begin{tabular}{|p{0.30\linewidth}|p{0.61\linewidth}|}
\hline 速度场连续 & 算$u_x+v_y+w_z$\\ \hline
有旋 & $\Omega_z=v_x-u_y$，$\omega_z=\Omega_z/2$\\ \hline
流线 & $dx/u=dy/v$，非定常先代$t=t_0$\\ \hline
势/流函数 & 先验无旋/不可压，再积分\\ \hline
小孔出流 & $V=\sqrt{2gH}$，$Q=\mu A\sqrt{2gH}$\\ \hline
皮托管 & $V=\sqrt{2(p_0-p)/\rho}$\\ \hline
U管差压 & $\Delta p=(\rho_m-\rho)g\Delta h$\\ \hline
维数分析 & $n$个量、$k$个基本量纲，得$n-k$个$\pi$数\\ \hline
\end{tabular}\renewcommand{\arraystretch}{1}

\sect{不会也要写的步骤分}
先写“该题属于……模型”，再写控制方程、边界条件、代入顺序。比如：“定常一维绝热等熵流，由$M$求滞止量，由面积比求另一截面$M$，再由等熵式求$p,T,\rho,V$。”这能拿框架分。
\sect{边界层积分方程模板}
无压梯度平板：$\tau_w=\rho U^2d\theta/dx$。若假定剖面含未知$\delta(x)$，先算$\theta=C\delta$，再由$\tau_w=\mu(\partial u/\partial y)_0$得到微分方程，积分求$\delta(x)$。
\sect{阻力系数换算}
局部$c_f=\tau_w/(0.5\rho U^2)$；平均$\bar C_f=D/(0.5\rho U^2A)$。层流局部$c_f=0.664/\sqrt{Re_x}$；平均$\bar C_f=1.328/\sqrt{Re_L}$。不要把局部当平均。
\sect{边界层题检查}
$x$用从前缘量的距离；$L$用整块板长度；$A$是湿表面积；双面板乘2。若题给临界$Re_{cr}$，先算转捩位置$x_{cr}=Re_{cr}\nu/U$。
\sect{最后10秒}
看见“平均时均速度”就是$\bar u$，不是中心最大速度；看见“壁面最大厚度”通常是$x=L$处$\delta$；看见“平板阻力”要用平均摩擦系数，不是尾部局部值。
\sect{外流阻力典型量}
球/圆柱阻力常写$D=C_D(0.5\rho U^2)A$，$A$为迎风面积。Stokes小球$Re<1$：$D=3\pi\mu dU$。高Re钝体以压差阻力为主，不能只算摩擦。
\sect{边界层分离答题句}
“由于逆压梯度，近壁流体动能不足，壁面速度梯度降为零并反向，边界层发生分离，尾流区压力降低，压差阻力增大。”
\sect{零知识保底路线}
题里有速度场：先连续/有旋/流线。题里有管路：先能量方程+损失。题里有Ma：先等熵总量，遇激波分段。题里有平板外流：先$Re$判层/湍，再摩擦系数。题里有应力张量：$P\bm n$。
\sect{查书优先级}
可带教材时，优先查：正/斜激波表、PM函数表、Moody图、面积--马赫数表、阻力系数图。速查表负责告诉你该查哪个表和查完怎么接公式。

\chap{Q1 静力学、运动学、张量题库速解}
\sect{压强单位换算}
$1{\rm kPa}=10^3{\rm Pa}$；$1{\rm mmHg}=133.3{\rm Pa}$；水头$h=p/(\rho g)$。表压$p_g=p_{\rm abs}-p_a$；绝压进入汽蚀、气体状态方程；表压进入多数水力计算。
\sect{平面闸门完整答案骨架}
1 取形心深度$h_C$。2 合力$F=\rho gh_CA$。3 压力中心：竖直平面$y_D=y_C+I_C/(y_CA)$；倾斜平面沿板长坐标同式，$h_C=y_C\sin\theta$。4 对铰点列力矩求拉力/反力。写明“压力分布为三角形或梯形”。
\sect{曲面闸门完整答案骨架}
水平分量$F_x$等于曲面在竖直投影面上的静水压力；竖直分量$F_y$等于曲面上方压力体水重，方向看压力体在曲面哪侧。合力$F=\sqrt{F_x^2+F_y^2}$，作用线由力矩或过圆心性质定。
\sect{相对平衡}
容器水平加速度$a$：自由液面满足$dp=0$，$g\,dz+a\,dx=0$，故$\tan\alpha=a/g$。竖直加速向上$a$：等效重力$g+a$；向下$a$：$g-a$。旋转容器：$p/\rho+gz-\omega^2r^2/2=C$，自由面$z=\omega^2r^2/(2g)+C$。
\sect{浮力和稳定}
浮力$F_B=\rho gV_{\rm disp}$，$V_{\rm disp}$为排开液体体积，作用于该体积形心。稳定判断：小倾角后浮心移动，重心$G$低于定倾中心$M$稳定；题若只问浮力，直接排水体积。
\sect{迹线/流线/脉线}
流线：固定时刻速度场切线，$dx/u=dy/v=dz/w$。迹线：单个质点轨迹，解$dx/dt=u(x,y,z,t)$。脉线：同一注入点不同时刻质点连线。定常流三者重合；非定常不能混写。
\sect{加速度题模板}
$a_x=u_t+uu_x+vu_y+wu_z$，$a_y=v_t+uv_x+vv_y+wv_z$，$a_z=w_t+uw_x+vw_y+ww_z$。若只给二维定常，去掉$t,z$项。题问“当地/迁移”分别对应$\partial/\partial t$和对流项。
\sect{应力张量标准解}
平面单位法向$\bm n=(n_x,n_y,n_z)^T$。应力矢量$\bm t=P\bm n$。法向应力$\sigma_n=\bm t\cdot\bm n$。切应力大小$\tau=\sqrt{|\bm t|^2-\sigma_n^2}$。夹角$\cos\alpha=\sigma_n/|\bm t|$。若问垂直平面分量，答案为$\sigma_n\bm n$。
\sect{应力对称原因}
无体偶力、微元角动量平衡给$p_{xy}=p_{yx}$，$p_{xz}=p_{zx}$，$p_{yz}=p_{zy}$。作用：张量独立分量从9个减到6个；做题时矩阵若不对称，通常题设有误或需取对称部分。
\sect{流体微团运动分解}
速度梯度$\nabla\bm V$可分解为平移、线变形率、角变形率、旋转。二维不可压时$u_x+v_y=0$；角速度$\omega_z=(v_x-u_y)/2$；角变形率$\varepsilon_{xy}=(u_y+v_x)/2$。
\sect{维数分析常用组合}
阻力问题$F=f(\rho,\mu,V,L)$：$F/(\rho V^2L^2)=\Phi(Re)$。自由液面重力波常有$Fr=V/\sqrt{gL}$。表面张力小尺度常有$We=\rho V^2L/\sigma$。可压缩高速常有$Ma=V/a$。

\chap{Q2 伯努利、动量、能量方程题库速解}
\sect{伯努利适用性}
理想无粘：沿流线；无旋势流：全场任意点。实际管路：用机械能方程加$h_f,\sum h_\zeta,H_p,H_T$。非定常势流：$\partial\varphi/\partial t+V^2/2+p/\rho+gz=C(t)$。
\sect{皮托管/文丘里/孔板}
皮托：滞止点$V=0$，$V=C\sqrt{2\Delta p/\rho}$。文丘里不可压：连续$A_1V_1=A_2V_2$，伯努利$p_1-p_2=\frac12\rho(V_2^2-V_1^2)$，解$Q=A_2\sqrt{\frac{2(p_1-p_2)/\rho}{1-(A_2/A_1)^2}}$，实际乘流量系数。
\sect{虹吸管}
最高点压强最低。列水面到最高点：$p_a/\rho g+z_0=p_c/\rho g+z_c+V^2/(2g)+h_{loss}$。要求$p_c>p_v$，否则断流/汽化。出口速度由总落差扣损失。
\sect{弯管受力}
对流体列动量，入口压力推入、出口压力推入控制体内部。若入口$x$向、出口转角$\alpha$：速度向量$\bm V_1=(V_1,0)$，$\bm V_2=(V_2\cos\alpha,V_2\sin\alpha)$。支座对流体力$\bm R$求出后，流体对管道力$-\bm R$。
\sect{射流冲击平板}
固定平板正撞：$F=\rho QV$。移动平板速度$u$：相对流量$\rho A(V-u)$，力$F=\rho A(V-u)^2$。功率$P=Fu$，最大功率常在$u=V/3$。斜板按法向动量变化。
\sect{喷嘴/消防水枪反力}
喷嘴出口表压近似0，入口压力加速流体：$p_1A_1-p_2A_2+R_x=\rho Q(V_2-V_1)$。若求握持力，取流体对喷嘴反力。先由伯努利或连续求$V_2,Q$。
\sect{水轮机/泵符号}
泵给流体能量：左侧加$H_p$。水轮机从流体取能：右侧加$H_T$。功率$P=\rho gQH$；若有效率，泵轴功率$P_{\rm shaft}=\rho gQH_p/\eta$，水轮机输出$P_{\rm out}=\eta\rho gQH_T$。
\sect{大水池到大水池}
两端大水面$V_1\approx V_2\approx0$，表压均0，则水位差$z_1-z_2=h_f+\sum h_\zeta-H_p+H_T$。这是长管、阀门、泵题最快写法。
\sect{压力表截面}
压力表给表压$p_g$，能量方程中可直接用$p_g/(\rho g)$，但汽蚀必须转绝压$p_{\rm abs}=p_g+p_a$。压力表位置速度不可忽略，面积变了要算$V=Q/A$。

\chap{Q3 第5/6章势流与理想不可压题库速解}
\sect{势函数流函数互求}
直角坐标：$u=\varphi_x=\psi_y$，$v=\varphi_y=-\psi_x$。若给$u,v$，先验无旋$v_x-u_y=0$再求$\varphi$；先验不可压$u_x+v_y=0$再求$\psi$。流线就是$\psi=C$。
\sect{极坐标速度关系}
$v_r=\partial\varphi/\partial r=(1/r)\partial\psi/\partial\theta$；$v_\theta=(1/r)\partial\varphi/\partial\theta=-\partial\psi/\partial r$。圆柱、点源、点涡题优先换极坐标。
\sect{基本流全表}
均匀流：$\varphi=Ur\cos\theta,\psi=Ur\sin\theta$。源：$\varphi=\frac Q{2\pi}\ln r,\psi=\frac Q{2\pi}\theta$。汇取负号。点涡：$\varphi=\frac\Gamma{2\pi}\theta,\psi=-\frac\Gamma{2\pi}\ln r$。偶极：$\varphi=M\cos\theta/r,\psi=-M\sin\theta/r$。
\sect{叠加题总流程}
1 读图选基元。2 写总$\psi$或$\varphi$。3 壁面/物体边界通常是$\psi=C$。4 驻点令$v_r=v_\theta=0$或$u=v=0$。5 压强用伯努利。6 合力用压力积分或库塔--儒可夫斯基。
\sect{半圆柱/半体绕流}
均匀流+点源构成Rankine半体。过驻点分界流线是物体边界。若是壁面半平面，把源/汇镜像后常等价于强度翻倍或分母$\pi$。
\sect{镜像法常见组合}
源靠固壁：同号源在镜像点。汇靠固壁：同号汇。点涡靠固壁：反号涡。双壁直角：连续镜像到四象限。目的只有一个：让壁面成为流线，法向速度为0。
\sect{圆柱绕流推导}
设$\varphi=(Ar+B/r)\cos\theta$。无穷远给$A=U$；壁面$r=a$不穿透：$v_r=\partial\varphi/\partial r=0$给$B=Ua^2$。所以$\varphi=U(r+a^2/r)\cos\theta$。
\sect{圆柱压力和阻力}
壁面$v_\theta=-2U\sin\theta$，$C_p=1-4\sin^2\theta$。前后对称，阻力为0即达朗贝尔佯谬。现实阻力来自粘性边界层分离，不是势流能算出的。
\sect{有环量圆柱/机翼}
叠加点涡后$v_\theta=-2U\sin\theta+\Gamma/(2\pi a)$。上下面速度不对称产生升力$L'=\rho U\Gamma$。库塔条件本质：后缘速度有限，确定环量。
\sect{边角绕流}
角域$0<\theta<\alpha$，壁面取$\psi=0$，解$\varphi=Ar^n\cos n\theta,\psi=Ar^n\sin n\theta$，$n=\pi/\alpha$。问速度奇异性用$V\sim r^{n-1}$。
\sect{势流压力题}
势流速度已知后，压力都是$p=p_0-\frac12\rho V^2$加常数。比较压强只需比较速度大小；速度大压强小。若点涡中心速度发散，理想模型中心不可直接用。

\chap{Q4 第7章可压缩、喷管、激波题库速解}
\sect{弹性模量/声速}
$E=-dp/(dV/V)=dp/(d\rho/\rho)$。水受压密度相对变化$\Delta\rho/\rho=\Delta p/E$。理想气体等熵声速$a=\sqrt{\gamma RT}$，马赫数$M=V/a$。
\sect{飞机声音/马赫锥}
$M>1$才有马赫锥，半角$\mu=\sin^{-1}(1/M)$。飞机到观察者最近水平距离或高度给出时，用几何关系$h/x=\tan\mu$或$x=h/\tan\mu$求听到位置/时间。
\sect{滞止量题}
若入口速度不可忽略，先算$M_1=V_1/\sqrt{\gamma RT_1}$，再$T_0=T_1(1+0.2M_1^2)$，$p_0=p_1(1+0.2M_1^2)^{3.5}$。若“入口速度可忽略”，直接$p_0=p_1,T_0=T_1$。
\sect{收缩喷管7.5型}
给容器$p_0,T_0$、背压$p_b$、出口面积$A_e$。若$p_b/p_0>0.528$未阻塞：$p_e=p_b$，由$p_0/p_e$求$M_e$，再算$\dot m$。若$\le0.528$阻塞：出口$M=1$，$p_e=p^*=0.528p_0$，$\dot m=0.0404A_ep_0/\sqrt{T_0}$。
\sect{Laval面积7.6型}
给$A_t=A^*$、出口状态，且无激波：出口若是超声速，就由等熵式求$M_e$，再用面积--马赫数求$A_e/A^*$。流量用喉部最大流量或出口$\rho_eV_eA_e$。
\sect{文丘里7.7型}
给最大流量和喉部最低压力：先入口截面由$\dot m=\rho VA$求$M_1$与$p_0,T_0$，再由$p_0/p_t$求$M_t$。只有$M_t=1$才临界；$M_t<1$时喉部不是阻塞。
\sect{正激波题7.9/7.10型}
已知$M_1,p_1,T_1$：用正激波关系求$M_2,p_2,T_2,\rho_2$；再用等熵式分别算$p_{01},p_{02}$。结论句：$T_0$守恒，$p_0$损失，熵增。
\sect{斜激波7.15/7.16/7.18型}
已知来流$M_1$和楔角$\theta$：查$\theta-\beta-M$得波角$\beta$；算$M_{n1}=M_1\sin\beta$；套正激波求法向波后；再$M_2=M_{n2}/\sin(\beta-\theta)$。压力比只用$M_{n1}$。
\sect{PM膨胀7.19型}
已知$M_1$和膨胀角$\delta$：$\nu(M_2)=\nu(M_1)+\delta$。已知出口压力低于/高于背压时，判断出口外压缩或膨胀；膨胀波内$p_0$不变。
\sect{管内正激波7.20型}
激波位置给$A_s/A_t$：波前按超声速支求$M_1$；过正激波求$M_2,p_{02}$；波后用新总压亚声速等熵到出口。若算出的出口压力等于背压，该激波位置成立。
\sect{Laval背压图文字版}
$p_b$很高：不流或全亚声速。降到临界：喉部$M=1$。再降：扩张段内正激波向出口移动。出口正激波后，再降：出口外斜激波。设计背压：全管等熵。更低：出口外膨胀波。
\sect{激波/膨胀选择}
超声速流被迫转向自身一侧，即流道变窄/凹角，压缩，用激波；超声速流道向外打开/凸角，膨胀，用PM扇。亚声速不能用这些关系。

\chap{Q5 第8章粘性流动、管流、边界层题库速解}
\sect{N--S简化四句话}
定常去$\partial/\partial t$；充分发展去$\partial/\partial x$；单向平行流只留一个速度分量；惯性项为零时方程变成“粘性扩散=压力梯度/重力”。写假设比直接算更重要。
\sect{平板间层流}
$\mu u''=dp/dx$。无压差Couette：$u=Uy/h$。无动板Poiseuille：$u=\frac{1}{2\mu}(-dp/dx)y(h-y)$。两者叠加：$u=Uy/h+\frac{1}{2\mu}(-dp/dx)y(h-y)$。
\sect{两层Couette书8.6型}
下层厚$h_1,\mu_1$，上层厚$h_2,\mu_2$，上板速度$U$。剪应力常数：$\tau=U/(h_1/\mu_1+h_2/\mu_2)$。速度分段线性，交界速度连续，剪应力连续。
\sect{倾斜平板8题型}
沿板方向无压差，重力分量驱动：$\mu u''+\rho g\sin\theta=0$。积分两次，用$u(0)=0,u(H)=-U_0,u(H/2)=0$求常数和$U_0$。剪应力$\tau=\mu u'$，上下壁代$y=0,H$。
\sect{圆管H--P}
速度抛物线，$u_{\max}=2\bar u$。$Q=\pi R^4\Delta p/(8\mu L)$，$\Delta p=32\mu L\bar u/d^2$，$\lambda=64/Re$。壁剪$\tau_w=8\mu\bar u/d$。
\sect{管道湍流壁律}
$u_\tau=\sqrt{\tau_w/\rho}$。线性底层$u^+=y^+$，缓冲层$5<y^+<30$，对数层$u^+=2.5\ln y^+ +5.5$。圆管平均$\bar u/u_\tau=2.5\ln(Ru_\tau/\nu)+1.75$。
\sect{管内边界层厚度题}
线性底层厚度$y_5=5\nu/u_\tau$；缓冲层外缘$y_{30}=30\nu/u_\tau$；完全湍流核心从$y^+>30$开始。若题问“离壁距离”，$y=R-r$。
\sect{管路损失综合题}
每段独立算$V_i,Q,A_i,Re_i,\lambda_i$。沿程$h_{fi}=\lambda_i(L_i/d_i)V_i^2/(2g)$。局部$h_\zeta=\zeta V^2/(2g)$，速度取元件所在管段。总损失相加进能量方程。
\sect{粗糙度使用}
层流不用粗糙度。湍流用$\varepsilon/d$查Moody或Colebrook。题给镀锌钢管粗糙度$h$时，通常$\varepsilon=h$，相对粗糙度$\varepsilon/d$。
\sect{泵安装高度}
水面到泵入口列能量：$0+0+z_0=p_B/\rho g+V^2/(2g)+z_B+h_f+\sum h_\zeta$。令$p_B=p_v+\Delta p_{\rm safe}$求最大$z_B-z_0$。越高越危险。
\sect{水轮机功率}
上游水面到下游截面：左边总水头减去损失和出口剩余水头，剩下就是$H_T$。$P=\rho gQH_T$。若给压力表，压力表截面还有速度头。
\sect{平板边界层}
层流$\delta=5x/\sqrt{Re_x}$，$\delta^*=1.72x/\sqrt{Re_x}$，$\theta=0.664x/\sqrt{Re_x}$，局部$c_f=0.664/\sqrt{Re_x}$，平均$\bar C_f=1.328/\sqrt{Re_L}$。
\sect{湍流/转捩平板}
全湍近似$\delta=0.37x/Re_x^{1/5}$，局部$c_f=0.0592/Re_x^{1/5}$，平均$\bar C_f=0.074/Re_L^{1/5}$。从前缘层流转捩用修正$\bar C_f=0.074/Re_L^{1/5}-1742/Re_L$。
\sect{平板阻力第10章题}
先算$Re_L=UL/\nu$，与$Re_{cr}$比。单面面积$Lb$，双面$2Lb$。总阻力$D=\bar C_f(0.5\rho U^2)A$。尾部厚度用$x=L$。
\sect{边界层动量积分}
给假设速度剖面，先算$\theta=C_\theta\delta$，壁剪$\tau_w=\mu U f'(0)/\delta$。代$\tau_w=\rho U^2d\theta/dx$得$\delta d\delta\sim(\nu/U)dx$，积分定常数。
\sect{分离/绕流阻力}
逆压梯度$dp/dx>0$导致近壁速度降低，$\tau_w=0$为分离点。圆柱、球分离后压差阻力大；流线型体摩擦阻力大。减阻常靠延迟分离或减小湿表面积。
\sect{Stokes小球}
$Re<1$：阻力$D=3\pi\mu dU$。沉降终速由重力减浮力等于阻力：$(\rho_s-\rho)g\pi d^3/6=3\pi\mu dU$，得$U=(\rho_s-\rho)gd^2/(18\mu)$。

\chap{Q6 课后题/往年题套写索引}
\sect{看见“书4.1--4.4”}
多为理想流体/伯努利/动量基础。先判断是否静止、是否沿流线、是否需控制体。能量题写伯努利，受力题写动量，压强题写静水压。
\sect{看见“书5.x 势流”}
先写$\varphi,\psi$基本流叠加。若题有圆柱/机翼，先圆柱公式；若题有墙，先镜像；若题有夹角，先$n=\pi/\alpha$；若题有压力，最后伯努利。
\sect{看见“第6章平面不可压无旋作业”}
速度场题：连续/有旋/势函数/流函数四连。剪切板题：不是势流，转到粘性N--S简化。题目混章节时，以“有没有粘性/有没有无旋”分流。
\sect{看见“第7章7.5--7.8”}
关键词是喷管、背压、阻塞、质量流量。先定$p_0,T_0$，再判断是否阻塞，最后算$M,p,T,\rho,V,A,\dot m$。0.528只在临界$M=1$用。
\sect{看见“第7章7.9--7.20”}
正激波、斜激波、膨胀波。先判断波型。正激波直接公式；斜激波先法向马赫数；膨胀波用PM函数；Laval内激波分波前、过波、波后三段。
\sect{看见“第8章8.4--8.10”}
多为层流解析解。写N--S简化，积分两次，代边界条件。两层流体必须写速度连续和剪应力连续。小球$Re<1$用Stokes。
\sect{看见“粘性流动作业-2前3题”}
第1边界层壁律：由$\bar u/u_\tau$方程求$u_\tau$，再算$y^+$厚度。第2泵吸水：水面到泵入口列能量，汽蚀压强约束。第3水轮机：水面到压力表截面列能量，分段损失。
\sect{看见“边界层第10章”}
先算$Re$，再选层流/湍流公式。问厚度用$\delta$；问阻力用$C_f$；问功率乘$U$；问分离写逆压梯度。
\sect{考场无法算表值时}
写“查表得$M=...,\lambda=...,\beta=...,\nu(M)=...$”，再把后续公式完整列出。老师通常给过程分；空着不写最亏。
\sect{最终卷面模板}
“本题属于\underline{\hspace{1cm}}模型。假设\underline{\hspace{1cm}}。控制方程为\underline{\hspace{1cm}}。由边界条件\underline{\hspace{1cm}}得未知量。代入数值得\underline{\hspace{1cm}}。物理判断：单位正确、方向正确、状态合理。”

\chap{R1 典型数值题答案骨架}
\sect{水的体积弹性模量}
水从1atm压到100atm：$\Delta p=99\times101325$Pa，$\Delta\rho/\rho=\Delta p/E$。$E=2.07\times10^9$Pa，故$\Delta\rho/\rho\approx4.85\times10^{-3}=0.485\%$。气体等温从1atm到2atm：$p/\rho=RT$，密度变为2倍，相对变化$100\%$。
\sect{超声速飞机听声}
$M=V/a=750/335=2.24$。马赫角$\mu=\sin^{-1}(1/M)\approx26.5^\circ$。飞机飞过观察者后，马赫线到达地面，水平距离$x=h/\tan\mu=h\sqrt{M^2-1}$。$h=1500$m，$x\approx3.0$km。
\sect{7.19管外膨胀数值型}
$A_e/A_t=2.005$查面积马赫数超声速支$M_e\approx2.20$。$p_e=p_0/(1+0.2M_e^2)^{3.5}$。若$p_e>p_b$，出口外膨胀。用$p_0/p_b$求最终$M_b\approx2.35$，$\alpha=\nu(M_b)-\nu(M_e)\approx3.9^\circ$。
\sect{第8章第1题管壁层}
直径$d=50$mm，$R=0.025$m，$\bar u=1.7$m/s，$\nu=1.006\times10^{-6}$。$Re=\bar u d/\nu$，湍流。由$\bar u/u_\tau=2.5\ln(Ru_\tau/\nu)+1.75$解$u_\tau$。线性底层$\delta_1=5\nu/u_\tau$，缓冲层到$\delta_2=30\nu/u_\tau$，核心从$R-\delta_2$到管心。
\sect{第8章第2题吸水泵}
$Q=50/3600$，$d=0.1$，$V=4Q/\pi d^2$，$Re=Vd/\nu$。相对粗糙度$\varepsilon/d=0.15{\rm mm}/100{\rm mm}=0.0015$查$\lambda$。水面到泵入口：$p_a/\rho g=p_B/\rho g+h+V^2/(2g)+(\lambda L/d+\zeta_{\rm valve}+\zeta_{\rm bend})V^2/(2g)$。$p_B=p_v+20{\rm kPa}$，解$h_{\max}$。
\sect{第8章第3题水轮机}
上游水面到下游压力表截面。上游水面$p=0,V=0,z=30$m；压力表$p=70$kPa表压，截面速度由$d=0.3$m算。斜管$d=0.2,L=30$，水平管$d=0.3,L=15$，分别算$V,Re,\lambda,h_f$。能量方程求$H_T$，$P=\rho gQH_T$。
\sect{平板边界层第8题}
$U=2.5$m/s，$L=0.5$m。空气$\rho=1.2,\nu=1.5\times10^{-5}$；水$\rho=998,\nu=1.005\times10^{-6}$。分别算$Re_L$与$Re_{cr}=5\times10^5$比较。层流用$\bar C_f=1.328/\sqrt{Re_L}$，若转捩用修正式。$D=\bar C_f(0.5\rho U^2)(Lb)$。
\sect{应力张量第4章第1题型}
矩阵$P$给定，单位法向$\bm n$给定。直接矩阵乘法$\bm t=P\bm n$。法向$\bm t_N=(\bm t\cdot\bm n)\bm n$；切向$\bm t_T=\bm t-\bm t_N$。若要“夹角”，$\alpha=\cos^{-1}[(\bm t\cdot\bm n)/|\bm t|]$。
\sect{圆柱绕流压力积分型}
$p(\theta)=p_\infty+\frac12\rho U^2(1-4\sin^2\theta)$。单位长度受力$d\bm F=-p(\cos\theta\bm i+\sin\theta\bm j)a\,d\theta$，积分$0$到$2\pi$。无环量$D=L=0$；有环量直接$L'=\rho U\Gamma$。

\chap{R2 教材可查表项目和速查表衔接}
\sect{Moody图怎么用}
横坐标$Re$，曲线参数$\varepsilon/d$，读Darcy摩阻系数$\lambda$。若题中公式用$f$需确认是否为Fanning系数；本课常用Darcy，$h_f=\lambda Ld^{-1}V^2/(2g)$。
\sect{面积--马赫数表怎么用}
输入$A/A^*$，会有亚声速和超声速两支。收缩喷管只能用亚声速支直到$M=1$；Laval喉后若设计超声速才用超声速支。题没说明时，结合背压和流动状态判支。
\sect{正激波表怎么用}
输入$M_1$，读$M_2,p_2/p_1,\rho_2/\rho_1,T_2/T_1,p_{02}/p_{01}$。表中$p_{02}/p_{01}$可省去手算总压损失；但波后继续等熵必须用$p_{02}$。
\sect{斜激波表怎么用}
输入$M_1$和偏转角$\theta$，通常有弱/强两解，工程取弱解。读$\beta$后用$M_{n1}=M_1\sin\beta$查正激波关系。若$\theta>\theta_{\max}$，无附体斜激波。
\sect{PM函数表怎么用}
输入$M$读$\nu(M)$。膨胀：$\nu_2=\nu_1+\delta$；压缩不能用PM膨胀表，而要斜激波。角度单位必须统一，表多为度。
\sect{阻力系数图怎么用}
圆柱/球外流阻力$D=C_D(0.5\rho U^2)A$，$A$是迎风面积，不是湿表面积。平板摩擦阻力用$C_f$，面积是湿表面积。两类不能混。

\chap{R3 高风险易混概念对照}
\sect{$p_0,p^*,p_e,p_b$}
$p_0$滞止压；$p^*$临界截面静压；$p_e$出口截面内静压；$p_b$背压/外界压。只有部分状态$p_e=p_b$。阻塞时收缩喷管$p_e=p^*$，不等于更低的$p_b$。
\sect{喉部、临界、阻塞}
喉部是最小面积截面。临界状态是该截面$M=1$。阻塞是继续降低背压也不能增大质量流量。喉部存在不等于一定临界；临界也不等于所有截面$M=1$。
\sect{亚声速/超声速面积效应}
亚声速：收缩加速、扩张减速。超声速：扩张加速、收缩减速。Laval喷管先收缩到喉部$M=1$，再扩张才能超声速。
\sect{激波和膨胀波}
激波是压缩、非等熵、总压损失；膨胀波是连续小扰动、近似等熵、总压不损失。激波可使超声速变亚声速；膨胀波使超声速更快。
\sect{层流/湍流/边界层}
层流是流动状态；边界层是近壁区域；湍流边界层不是“管内湍流核心”。管内H--P只能用于层流充分发展，不能套到湍流。
\sect{平均速度/最大速度/摩擦速度}
$\bar u$是截面平均；$u_{\max}$是中心最大；$u_\tau=\sqrt{\tau_w/\rho}$不是实际流速，而是由壁剪定义的速度尺度。壁律题用$u_\tau$。
\sect{表面积/迎风面积/截面积}
管流连续用截面积$\pi d^2/4$；平板摩擦用湿表面积$Lb$或$2Lb$；绕流阻力用迎风面积；静水压力用受压面积投影或实际面积，别混。
\sect{局部损失/沿程损失}
沿程损失跟$L/d$有关；局部损失跟$\zeta$有关。弯头、阀门、入口、出口不是沿程损失。长管可以局部忽略，短管不能。
\sect{绝压/表压}
气体状态方程、汽化压、声速问题用绝压/绝温。水力能量方程两端同大气时可用表压。压力表读数通常是表压。

\chap{R4 按题干关键词秒选公式}
\sect{“不可压无旋”}
写$\nabla^2\varphi=0,\nabla^2\psi=0$，用$\varphi,\psi$速度关系，压力用伯努利。若题说“有环量”，加点涡。
\sect{“定常一维绝热”}
若无摩擦激波，等熵；若有激波，分段等熵+激波关系；若有摩擦热交换，本表不套等熵，按题给关系。
\sect{“平均时均速度”}
通常湍流统计量$\bar u$；若在管内，用壁律或平均速度公式，不是瞬时速度。雷诺应力$\overline{u'v'}$来自脉动相关。
\sect{“小球Re<1”}
直接Stokes阻力；若问沉降终速，重力-浮力=阻力；若问阻力系数，$C_D=24/Re$。
\sect{“完全发展”}
速度分布不随流向变，$\partial u/\partial x=0$，但压强仍沿流向下降。充分发展层流可解析；入口段不能直接用H--P。
\sect{“无压差Couette”}
速度线性，剪应力常数。若有压力梯度，速度为线性+抛物线。若两层流体，剪应力在两层相等。
\sect{“逆压梯度”}
$dp/dx>0$，流体沿流动方向压力升高，近壁减速，可能分离。顺压梯度$dp/dx<0$则加速，分离不易发生。
\sect{“背压降低”}
喷管题中背压降低不是静压处处降低。未阻塞时影响全管；阻塞后质量流量固定，影响下游激波位置/管外波系。
\sect{“压力中心/形心”}
形心用于算合力大小；压力中心用于合力作用点。压力中心在形心以下，除非面积水平或压强均匀。

\chap{R5 计算细节防错}
\sect{单位统一}
面积：$1{\rm cm^2}=10^{-4}{\rm m^2}$；粗糙度mm转m；运动粘度$1{\rm mm^2/s}=10^{-6}{\rm m^2/s}$；温度用K；角度查表通常用度，计算三角函数确认模式。
\sect{截面积}
圆管$A=\pi d^2/4$，半径题$A=\pi R^2$。环形面积$A=\pi(R_o^2-R_i^2)$。平板流每单位宽度时$Q'=\int udy$，总流量$Q=bQ'$。
\sect{水头与压力}
$p/\rho g$单位是m。能量方程每一项都用m时，泵扬程$H_p$也是m；功率再乘$\rho gQ$。用压力形式时每项是Pa，损失为$\Delta p=\rho gh$。
\sect{查表插值}
考试可粗插：若$x$在$x_1,x_2$之间，$y\approx y_1+(y_2-y_1)(x-x_1)/(x_2-x_1)$。写“线性插值”即可。
\sect{答案量纲检查}
流量$m^3/s$，质量流量$kg/s$，压强Pa，剪应力Pa，功率W，扬程m，环量$m^2/s$，势函数$m^2/s$，流函数$m^2/s$每单位宽度流量。
\sect{方向检查}
力的方向若不确定，先假设正向，算出负值表示反向。动量方程一定是“外力合力=流出动量流率-流入动量流率”。
\sect{结果合理性}
收缩喷管未阻塞$M<1$；正激波后$M_2<1$；膨胀波后$M$增大、$p$降低；泵入口高度越大入口压越低；边界层厚度沿程增大。

\chap{R6 可直接抄的结论句}
\sect{势流结论句}
“由无旋和不可压知势函数与流函数均满足Laplace方程，且线性方程允许叠加。固体边界为流线，故总流函数在边界上为常数。”
\sect{达朗贝尔佯谬句}
“理想不可压无粘定常绕流中，圆柱前后压力分布对称，积分阻力为零；实际阻力来自粘性边界层分离。”
\sect{边界层句}
“高雷诺数外流可分为外部近似无粘区和近壁粘性边界层；壁面无滑移使速度从0过渡到外流速度。”
\sect{管流损失句}
“沿程损失反映壁面摩擦，局部损失反映流道突变、阀门、弯头等引起的附加耗能；二者均以速度头表示。”
\sect{激波句}
“激波是超声速压缩扰动的有限强度间断，满足守恒方程但熵增，因此总压降低；正激波后必为亚声速。”
\sect{膨胀波句}
“膨胀波由连续弱扰动组成，可近似看作等熵过程；超声速气流通过膨胀扇后马赫数增大，静压和静温降低。”
\sect{阻塞句}
“阻塞的本质是最小面积处达到声速，信息不能向上游传递，继续降低背压不能增加质量流量。”
\sect{湍流句}
“湍流时速度可分解为时均量和脉动量，脉动相关产生雷诺应力，工程上常通过经验摩擦系数或壁律求解。”

\chap{G1 题目示意图库：看图选模型}
\sect{U管差压/振荡}
\begin{center}\begin{tikzpicture}[scale=0.62,>=Latex,every node/.style={font=\tiny}]
\draw[thick] (0,1.5)--(0,0)--(0.8,0)--(0.8,1.0);
\draw[thick] (1.4,1.2)--(1.4,0)--(0.9,0)--(0.9,0.8);
\draw[blue] (0.02,1.15)--(0.78,1.15) (0.92,0.68)--(1.38,0.68);
\draw[<->] (1.65,0.68)--(1.65,1.15); \node[right] at (1.65,0.92){$\Delta h$};
\node at (0.75,-0.25){$\Delta p=(\rho_m-\rho)g\Delta h$；振荡$T=2\pi\sqrt{L/2g}$};
\end{tikzpicture}\end{center}
\sect{倾斜闸门}
\begin{center}\begin{tikzpicture}[scale=0.62,>=Latex,every node/.style={font=\tiny}]
\draw[blue] (-0.2,1.2)--(2.1,1.2); \node[left] at (-0.2,1.2){液面};
\draw[thick] (0.55,0.05)--(1.65,0.65);
\draw[->,red] (1.1,0.35)--(0.75,0.98); \node[right] at (1.15,0.7){$F$};
\draw[dashed] (1.1,0.35)--(1.1,1.2); \node[right] at (1.1,0.78){$h_C$};
\node at (1.0,-0.25){$F=\rho gh_CA,\;y_D=y_C+I_C/(y_CA)$};
\end{tikzpicture}\end{center}
\sect{曲面受压}
\begin{center}\begin{tikzpicture}[scale=0.62,>=Latex,every node/.style={font=\tiny}]
\draw[blue] (-0.1,1.25)--(2.6,1.25);
\draw[thick] (0.5,0.2) arc (210:330:0.8);
\draw[->,red] (1.2,0.4)--(1.8,0.4); \node at (1.8,0.6){$F_x$};
\draw[->,red] (1.2,0.4)--(1.2,-0.25); \node[right] at (1.2,0.0){$F_y$};
\node at (1.2,-0.55){$F_x=$投影面压力；$F_y=$压力体重};
\end{tikzpicture}\end{center}
\sect{弯管动量}
\begin{center}\begin{tikzpicture}[scale=0.62,>=Latex,every node/.style={font=\tiny}]
\draw[thick] (0,0)--(1.1,0)--(1.8,0.65);
\draw[thick] (0,0.25)--(1.0,0.25)--(1.63,0.85);
\draw[->,blue] (-0.2,0.13)--(0.55,0.13); \node at (0.25,0.42){$V_1$};
\draw[->,blue] (1.45,0.48)--(2.05,1.05); \node at (2.0,0.72){$V_2$};
\node at (1.0,-0.32){$\sum\bm F=\rho Q(\bm V_2-\bm V_1)$，管受力取反};
\end{tikzpicture}\end{center}
\sect{射流冲板}
\begin{center}\begin{tikzpicture}[scale=0.62,>=Latex,every node/.style={font=\tiny}]
\draw[->,blue] (0,0)--(1.1,0); \node at (0.55,0.25){$\rho Q,V$};
\draw[thick] (1.35,-0.55)--(1.35,0.55);
\draw[->,red] (1.35,0)--(2.05,0); \node at (1.8,0.25){$F$};
\node at (1.1,-0.82){固定正撞$F=\rho QV$；动板$F=\rho A(V-u)^2$};
\end{tikzpicture}\end{center}
\sect{点源+均匀流半体}
\begin{center}\begin{tikzpicture}[scale=0.62,>=Latex,every node/.style={font=\tiny}]
\foreach \y in {-0.4,0,0.4}{\draw[->,blue] (-0.4,\y)--(0.45,\y);}
\draw[fill=black] (0.8,0) circle (0.04); \node[above] at (0.8,0){源};
\draw[thick] (0.0,0).. controls (0.65,0.65) and (1.7,0.75)..(2.4,0.55);
\draw[thick] (0.0,0).. controls (0.65,-0.65) and (1.7,-0.75)..(2.4,-0.55);
\node at (1.1,-1.0){驻点$r=Q/(2\pi U)$，分界$\psi=Q/2$};
\end{tikzpicture}\end{center}
\sect{圆柱绕流/有环量}
\begin{center}\begin{tikzpicture}[scale=0.62,>=Latex,every node/.style={font=\tiny}]
\foreach \y in {-0.5,0,0.5}{\draw[->,blue] (-0.6,\y)--(0.25,\y);}
\draw[thick] (0.9,0) circle (0.45);
\draw[->,red] (1.28,0.2) arc (25:140:0.45);
\node at (1.0,-0.75){$C_p=1-4\sin^2\theta$；$L'=\rho U\Gamma$};
\end{tikzpicture}\end{center}
\sect{Laval喷管/内激波}
\begin{center}\begin{tikzpicture}[scale=0.62,>=Latex,every node/.style={font=\tiny}]
\draw[thick] (0,0.55).. controls (0.7,0.45) and (0.9,0.2)..(1.45,0.2).. controls (2.1,0.2) and (2.25,0.55)..(3.2,0.7);
\draw[thick] (0,-0.55).. controls (0.7,-0.45) and (0.9,-0.2)..(1.45,-0.2).. controls (2.1,-0.2) and (2.25,-0.55)..(3.2,-0.7);
\draw[dashed] (1.45,-0.45)--(1.45,0.45); \node at (1.45,-0.7){喉};
\draw[red,thick] (2.35,-0.55)--(2.35,0.55); \node[red] at (2.55,0.65){激波};
\node at (1.7,-1.0){波前等熵$\to$正激波$\to$波后新$p_{02}$等熵};
\end{tikzpicture}\end{center}
\sect{斜激波楔角}
\begin{center}\begin{tikzpicture}[scale=0.62,>=Latex,every node/.style={font=\tiny}]
\draw[->,blue] (-0.1,0.45)--(0.9,0.45); \node at (0.35,0.72){$M_1$};
\draw[thick] (1,0)--(2.6,0); \draw[thick] (1,0)--(2.6,0.6);
\draw[red,thick] (1,0)--(1.65,1.05); \node at (1.55,0.2){$\theta$}; \node[red] at (1.32,0.92){$\beta$};
\node at (1.35,-0.35){先$M_{n1}=M_1\sin\beta$};
\end{tikzpicture}\end{center}
\sect{膨胀扇}
\begin{center}\begin{tikzpicture}[scale=0.62,>=Latex,every node/.style={font=\tiny}]
\draw[thick] (0,0)--(1.4,0); \draw[thick] (1.4,0)--(2.8,-0.55);
\foreach \a in {0.0,0.18,0.36,0.54,0.72}{\draw[red] (1.4,0)--(2.15,-\a);}
\draw[->,blue] (0.2,0.28)--(1.0,0.28); \node at (2.0,-0.95){$\nu_2=\nu_1+\delta$，$p_0$不变};
\end{tikzpicture}\end{center}
\sect{泵吸水高度}
\begin{center}\begin{tikzpicture}[scale=0.62,>=Latex,every node/.style={font=\tiny}]
\draw[blue] (0,0)--(1.0,0); \draw[thick] (0.5,0)--(0.5,1.2)--(2.0,1.2);
\node[circle,draw,inner sep=1pt] at (2.2,1.2){P};
\draw[<->] (2.55,0)--(2.55,1.2); \node[right] at (2.55,0.6){$h$};
\node at (1.4,-0.35){水面$\to$泵入口，$p_B\ge p_v+\Delta p$};
\end{tikzpicture}\end{center}
\sect{水轮机管路}
\begin{center}\begin{tikzpicture}[scale=0.62,>=Latex,every node/.style={font=\tiny}]
\draw[blue] (0,1.2)--(0.8,1.2); \draw[thick] (0.4,1.2)--(1.2,0.55)--(2.2,0.55);
\node[circle,draw,inner sep=1pt] at (1.2,0.55){T};
\draw[->,blue] (2.0,0.55)--(2.65,0.55);
\node at (1.35,0.05){$H_T$由能量方程；$P=\rho gQH_T$};
\end{tikzpicture}\end{center}
\sect{平板边界层}
\begin{center}\begin{tikzpicture}[scale=0.62,>=Latex,every node/.style={font=\tiny}]
\draw[thick] (0,0)--(3.0,0); \draw[->,blue] (-0.2,0.55)--(0.65,0.55);
\draw[red,thick] (0,0).. controls (1.0,0.25) and (2.0,0.5)..(3.0,0.8);
\node at (1.5,-0.25){层流$\delta=5x/\sqrt{Re_x}$；阻力用$\bar C_f$};
\end{tikzpicture}\end{center}
\sect{两层Couette}
\begin{center}\begin{tikzpicture}[scale=0.62,>=Latex,every node/.style={font=\tiny}]
\draw[thick] (0,0)--(2.6,0); \draw[thick] (0,1.2)--(2.6,1.2);
\draw[dashed] (0,0.55)--(2.6,0.55); \node[left] at (0,0.28){$\mu_1$}; \node[left] at (0,0.9){$\mu_2$};
\draw[->,red] (0.3,1.2)--(1.4,1.2); \node at (1.55,1.38){$U$};
\node at (1.3,-0.28){速度连续、剪应力连续；$\tau=U/(h_1/\mu_1+h_2/\mu_2)$};
\end{tikzpicture}\end{center}

\chap{G2 按图像补充的题型第一步}
\sect{玻璃弯管放水}
看见竖直段+倾斜段+阀门：未开阀时静水压分布$p=p_a+\rho g\Delta z$，只看竖直高度。开阀排空时间：沿流线伯努利求出口$V=\sqrt{2gH_{\rm eff}}$，再用液面下降连续方程积分。
\sect{半圆柱/圆柱浮起}
无穷远均匀流绕圆柱或半圆柱，用圆柱绕流速度分布求表面压力，再对压力积分求升力。若有重力和浮力，临界浮起条件：流体动力升力+浮力=重力。
\sect{点源靠壁最大速度}
壁面上点源/点汇用镜像。壁面速度由总势函数求导，令$dv/dx=0$找最大速度。给最大速度反求距离$a$或源强$q$。
\sect{水下半圆球/半圆柱开孔}
轴对称势流用给定流函数/势函数，开孔使压力等于外界时，先求表面压力分布$p=p_\infty+\frac12\rho(U^2-V_s^2)$，令$p=p_{\rm in}$解角度。
\sect{粘性圆管入口发展}
若题问“完全发展前的边界层厚度”，用边界层概念；若题问“完全发展层流压降”，用H--P。入口段不能直接套抛物线速度分布。
\sect{雷诺输运定理题}
系统量变化=控制体内变化+穿过控制面的净流出。定常时控制体内变化为0。动量方程、质量守恒都是它的特例。
\sect{LES/DNS/RANS概念}
DNS解全尺度N--S，最贵；LES解析大涡、模型小涡；RANS对时均方程建模雷诺应力，工程最常用。考试多为概念判断。
\sect{雷诺应力}
$-\rho\overline{u_i'u_j'}$来自湍流脉动动量输运，不是分子粘性应力。近壁湍流总剪应力=分子粘性剪应力+雷诺剪应力。
\sect{混合长度}
Prandtl混合长度假设$\tau_t=\rho l^2|du/dy|du/dy$，近壁$l=\kappa y$，可推对数律。只需知道它是湍流半经验模型。
\sect{边界层厚度种类}
名义厚度$\delta$：$u=0.99U$处；位移厚度$\delta^*$：外流被排挤的等效厚度；动量厚度$\theta$：动量亏损厚度。积分题按定义算。

\chap{T1 原题压缩收录：第4章动力学/应力}
\sect{应力张量+指定平面}
原题型：给$P$和单位法向$\bm n$，求应力矢量、垂直分量、夹角。答案：$\bm p_n=P\bm n$，$\bm p_N=(\bm p_n\cdot\bm n)\bm n$，$\bm p_T=\bm p_n-\bm p_N$，$\alpha=\cos^{-1}\frac{\bm p_n\cdot\bm n}{|\bm p_n|}$。
\sect{圆柱切平面M点应力}
圆柱面$y^2+z^2=4$在点$M(2,1,\sqrt3)$，法向取$\nabla(y^2+z^2-4)=(0,2y,2z)$归一化。代入坐标先求$P(M)$，再$\bm p=P\bm n$。
\sect{给应力分布求平面上一点}
先把点坐标代入$p_{xx},p_{xy},p_{yy}$得到矩阵$P$。平面$x+3y+z+1=0$法向$\bm n=(1,3,1)/\sqrt{11}$。同样$P\bm n$，再分解法向/切向。
\sect{由$\nabla\cdot P=-\rho\bm f$求体力}
散度按行求：$(\nabla\cdot P)_x=\partial p_{xx}/\partial x+\partial p_{xy}/\partial y+\partial p_{xz}/\partial z$，其余同理。最后$\bm f=-(\nabla\cdot P)/\rho$。
\sect{正交曲线坐标散度}
若选做，核心是尺度因子$h_i$。圆柱坐标向量散度：$\nabla\cdot\bm A=\frac1r\frac{\partial(rA_r)}{\partial r}+\frac1r\frac{\partial A_\theta}{\partial\theta}+\frac{\partial A_z}{\partial z}$。张量散度需查教材公式，不要现场硬推。

\chap{T2 原题压缩收录：第5章理想流体/第6章势流}
\sect{环形路径速度环量}
定常平面流，若给$v_r=ar,v_\theta=0$，沿半圆/环形路径的速度环量$\Gamma=\int\bm V\cdot d\bm l$。径向段$d\bm l=dr\bm e_r$有贡献；圆弧段$d\bm l=rd\theta\bm e_\theta$无贡献。按箭头方向分段积分$ar\,dr$。
\sect{弯玻璃管放水}
未开阀压力分布只由高度决定：从自由面向下$p=p_a+\rho gh$；倾斜段同样用竖直高度，不用管长。排空时间：液面高度$H(t)$，出口$V=\sqrt{2gH_{\rm eff}}$，$A_{\rm tank}dH/dt=-A_oV$积分。
\sect{速度场$u=3x^2-3y^2,v=-6xy$}
速度势：$\varphi=x^3-3xy^2+C$。流函数：$\psi=3x^2y-y^3+C$。速度势线$\varphi=C$，流线$\psi=C$。
\sect{$u=2cxy,v=c(a^2+x^2-y^2)$}
连续：$u_x+v_y=2cy-2cy=0$。旋度：$v_x-u_y=2cx-2cx=0$，无旋。$\varphi=\int2cxy\,dx=cx^2y+f(y)$，由$\varphi_y=cx^2+f'=c(a^2+x^2-y^2)$得$f=ca^2y-cy^3/3$。$\psi_y=u$求$\psi=cxy^2+g(x)$，再由$-\psi_x=v$定$g=-ca^2x-cx^3/3$。
\sect{平板Couette流量}
下板固定，上板$U_0$，间距$h$，无压差：$u=U_0y/h$。单位宽流量$q=\int_0^h udy=U_0h/2$；总流量$Q=bU_0h/2$。若有压差，叠加抛物线项。
\sect{势流方程推导}
不可压$\nabla\cdot\bm V=0$且无旋$\bm V=\nabla\varphi$，故$\nabla^2\varphi=0$。二维不可压定义$\psi$后，同样$\nabla^2\psi=-\omega_z$，无旋时$\nabla^2\psi=0$。
\sect{均匀流绕圆柱推导题}
写通解$(Ar+B/r)\cos\theta$。无穷远$A=U$；壁面$r=a$不穿透给$B=Ua^2$。速度和压力见P2。若问流线图，$\psi=U(r-a^2/r)\sin\theta=C$。
\sect{点源/点汇/点涡叠加题}
先定强度符号：源$+Q$，汇$-Q$，逆时针涡$+\Gamma$。总流函数相加。驻点令速度为0；边界是过驻点的流线。墙面问题先镜像再叠加。
\sect{边角流动题}
夹角$\alpha$，$n=\pi/\alpha$。$\varphi=Ar^n\cos n\theta$，$\psi=Ar^n\sin n\theta$。速度$V=nAr^{n-1}$。若$\alpha=120^\circ=2\pi/3$，$n=3/2$。

\chap{T3 原题压缩收录：第6章作业-2势流应用}
\sect{均匀流绕长钝体上坡}
无穷远均匀流$U$，已知A、B在同一流线。理想不可压无旋沿同一流线伯努利：$p_A+\rho gz_A+\frac12\rho V_A^2=p_B+\rho gz_B+\frac12\rho V_B^2$。若B点速度由势流叠加求，代入求高度或压差。
\sect{半圆柱屋顶受升力}
绕半圆柱用镜像等价全圆柱绕流的一半。表面速度$V_\theta=2U_\infty\sin\theta$。压力$p=p_\infty+\frac12\rho(U_\infty^2-V_\theta^2)$。升力为压力竖直分量积分；若有开孔使内外压相等，令表面$p(\theta)=p_{\rm in}$解孔位。
\sect{半圆柱水底浮起}
水深环境加静压，但动力压差由绕流速度给。半圆柱受力=静水浮力+动压积分。临界$U_\infty$由上浮力=重力求。若题说内部压强与上压强一样，内外静压抵消，只剩动力压。
\sect{壁面点源最大速度}
点源距壁$a$，镜像同号。壁面切向速度由两个源叠加得到，常形如$u(x)=q x/[\pi(x^2+a^2)]$，令$du/dx=0$得$x=a$处最大，$u_{\max}=q/(2\pi a)$或按题定义$q/(2\pi)$修正。

\chap{T4 原题压缩收录：第7章作业}
\sect{7.5收缩喷管}
题面：等熵流入容器，喷管入口速度忽略，给入口$p_0,T_0$、容器压力$p_b$、出口面积$A_e$，求流量和阻塞流量。答案骨架：算$p_b/p_0$；未阻塞则$p_e=p_b$反求$M_e$；阻塞则$M_e=1$，$\dot m_{\max}=0.0404A_ep_0/\sqrt{T_0}$。
\sect{7.6 Laval无激波}
题面常给$A_t$和出口$T_e$或$p_e$。若由等熵式得$M_e>1$且题说无激波，则喉部临界$A_t=A^*$。出口面积$A_e=A^*(A/A^*)_{M_e}$。流量用喉部最大流量。
\sect{7.7文丘里流量计}
题面给最大流量、入口截面、最低喉压。先入口状态求$p_0,T_0$，再由最低喉压求$M_t$。若$M_t<1$，不是临界；喉部面积$A_t=\dot m/(\rho_tV_t)$。
\sect{7.8类似喷管判断}
若题给背压变化，画$p_b$降低流程：未阻塞$\to$临界$\to$激波位置移动$\to$设计/欠膨胀。题问“是否阻塞/是否激波”，不要一上来代0.528。
\sect{7.9/7.10正激波}
已知$M_1$，查正激波表或公式。波后$M_2<1$；$p_2/p_1>1,T_2/T_1>1,\rho_2/\rho_1>1$；$p_{02}<p_{01}$。若题问声速/速度，先$T_2$再$a_2=\sqrt{\gamma RT_2}$，$V_2=M_2a_2$。
\sect{运动正激波撞墙反射}
坐标系固结在激波上：来流相对激波超声速。先过入射正激波求波后速度；反射波再以墙面速度为边界，使最终气体速度为0。过程复杂时写两个正激波关系+速度转换。
\sect{7.15/7.16斜激波}
楔角题：$M_1,\theta$查弱激波$\beta$。$M_{n1}=M_1\sin\beta$，用正激波表得$p_2/p_1,M_{n2}$，$M_2=M_{n2}/\sin(\beta-\theta)$。若问强弱解，通常弱解物理可实现。
\sect{7.17选做方程组}
若要求列方程：质量$\rho_1V_{n1}=\rho_2V_{n2}$；切向速度$V_{t1}=V_{t2}$；法向动量$p_1+\rho_1V_{n1}^2=p_2+\rho_2V_{n2}^2$；能量$h_1+V_1^2/2=h_2+V_2^2/2$。
\sect{7.18脱体激波}
给大偏转角或钝头体，若$\theta>\theta_{\max}$，斜激波不能附着，形成脱体弓形激波。近前缘可用正激波近似，远处转为斜激波。
\sect{7.19管外膨胀}
面积比$A_e/A_t$得$M_e$；等熵得$p_e$。若$p_e>p_b$欠膨胀? 注意：出口压高于环境，射流向外膨胀，管外膨胀波；若$p_e<p_b$，外界压缩，出现斜激波。最终角度用PM函数差。
\sect{7.20管内激波}
给激波截面$A_s$，波前用$A_s/A_t$超声速支求$M_1$；过正激波求$M_2,p_{02}/p_{01}$；波后把$A_s$视作新亚声速支上的截面，求到出口压力。
\sect{7.21风洞题}
题面缺图时：超声速风洞内通常用喷管/试验段面积比求$M$，再用等熵关系求$p,T,\rho,V$。若有激波，按位置分段。答案必须先补完整题图或数据。
\sect{作业3第4题U型压差计+正激波}
收缩扩张喷管出口有正激波，U型压差计连接喉部和下游大容器。压差计给$p_t-p_b=\rho_{\rm Hg}gh$或考虑水柱修正。若下游$p_b=100$kPa，先判断是否有激波；有则位置可能在管内/管外，按背压匹配。

\chap{T5 原题压缩收录：第8章粘性作业}
\sect{Re判层/湍}
水管$d=0.2,V=1,\nu=10^{-6}$：$Re=2\times10^5$，湍流。油管$d=0.15,V=0.2,\nu=2\times10^{-5}$：$Re=1500$，层流。圆管临界约2300。
\sect{书8.4类：圆管层流压降}
给$\rho,\mu,d,V,L$，先$Re$判层流。$\Delta p=\lambda(L/d)\rho V^2/2$，层流$\lambda=64/Re$。也可H--P：$\Delta p=32\mu LV/d^2$。
\sect{书8.5小球Re<1}
Stokes阻力$F=3\pi\mu dU$。若求沉降速度：$(\rho_s-\rho)g\pi d^3/6=3\pi\mu dU$。若求阻力系数：$C_D=24/Re$。
\sect{书8.6两层流体}
总速度差$U=\tau(h_1/\mu_1+h_2/\mu_2)$。所以$\tau=U/(h_1/\mu_1+h_2/\mu_2)$。分段速度：$u_1=\tau y/\mu_1$；$u_2=u_1(h_1)+\tau(y-h_1)/\mu_2$。
\sect{书8.10小球Re<1}
通常仍用Stokes。若问阻力或终速，先算$Re=\rho Ud/\mu$核验$<1$，再套。若结果导致$Re$不小，说明假设不自洽。
\sect{管内层流三方向}
水平管：$dp/dx=-32\mu V/d^2$。向上竖直：需要额外克服重力，$-dp/dz=32\mu V/d^2+\rho g$。向下竖直：$-dp/dz=32\mu V/d^2-\rho g$，可能压力沿流向升高。
\sect{倾斜平板层流}
方程$\mu u''+\rho g\sin\theta=0$，通解$u=-\rho g\sin\theta\,y^2/(2\mu)+C_1y+C_2$。代边界和$u(H/2)=0$，求上板速度和壁面剪应力。
\sect{物面条件证明}
静止固壁无滑移$\bm V=0$。涡通量提示：$\bm\Omega\cdot\bm n=0$来自壁面速度为零且不可穿透，考试写物理解释：固壁上速度无切向变化导致法向涡量为零。
\sect{流函数Poisson方程}
二维不可压：$u=\psi_y,v=-\psi_x$。涡量$\Omega=v_x-u_y=-(\psi_{xx}+\psi_{yy})=-\nabla^2\psi$，故$\nabla^2\psi=-\Omega$。
\sect{平板边界层混合}
若前段层流后段湍流，总阻力不是简单全湍。常用修正平均摩擦系数$\bar C_f=0.074/Re_L^{1/5}-1742/Re_L$，适用于临界$5\times10^5$。
\sect{圆管湍流平均速度公式}
$\bar u/u_\tau=2.5\ln(Ru_\tau/\nu)+1.75$用于题给平均速度反解摩擦速度。解出$u_\tau$后$\tau_w=\rho u_\tau^2$，$\lambda=8(u_\tau/\bar u)^2$。

\chap{T6 再补一页：考试常见问法到答案}
\sect{“证明满足连续方程”}
直接算$\nabla\cdot\bm V$。若为二维：$u_x+v_y=0$。写“故不可压连续方程满足”。若不为0，说明该速度场不可能表示不可压流。
\sect{“判断有旋否”}
二维：$\Omega_z=v_x-u_y$。三维：$\nabla\times\bm V$。若为0，无旋，可求势函数；若非0，不能全场定义单值势函数。
\sect{“写流线方程”}
定常：$dy/dx=v/u$。分离变量积分。非定常问某时刻流线，先把$t=t_0$代入速度场再积分。
\sect{“写迹线方程”}
解ODE：$dx/dt=u(x,y,t)$，$dy/dt=v(x,y,t)$，带初始条件$x(t_0)=x_0,y(t_0)=y_0$。不要用$dy/dx=v/u$替代。
\sect{“求流量”}
均匀速度$Q=VA$；非均匀$Q=\int_Au\,dA$。圆管轴对称$Q=2\pi\int_0^Ru(r)rdr$。平板单位宽$q=\int_0^hu(y)dy$。
\sect{“求剪应力”}
牛顿流体简单剪切$\tau=\mu du/dy$。圆管壁剪可由速度梯度或受力平衡。湍流总剪应力还含雷诺应力，题若是壁面，壁面脉动速度为0，壁剪仍由粘性梯度给。
\sect{“求压强分布”}
静止流体$p=p_a+\rho gh$。势流定常$p=p_\infty+\frac12\rho(U_\infty^2-V^2)$。管路沿程压降由损失或H--P。可压缩等熵由$p_0/p(M)$。
\sect{“求升力/阻力”}
压力积分：$\bm F=-\int p\bm n\,dA+\int\bm\tau\,dA$。势流圆柱无阻力；有环量升力$\rho U\Gamma$。工程外流可用$C_D,C_L$。
\sect{“求功率”}
泵/水轮机$P=\rho gQH$；阻力功率$P=DU$；剪切拖动功率$P=FU$；射流对动板功率$P=Fu$。注意效率。
\sect{“判断能不能用伯努利”}
同一流线、定常、不可压、无粘或损失可忽略。若有泵/损失，用机械能方程。若跨激波，不可用同一个伯努利常数。
\sect{“判断阻塞”}
只看最小面积截面是否达到$M=1$。收缩管用$p_b/p_0\le0.528$判断；Laval要结合背压状态和面积比。阻塞后$\dot m$不随背压继续增大。
\sect{“判断激波位置”}
给背压时，假设激波位置，算出口压力，与背压匹配。背压高，激波靠上游；背压降低，激波向出口移动；再低则到管外。
\sect{“判断边界层层流湍流”}
平板用$Re_x=Ux/\nu$，圆管用$Re=Vd/\nu$。平板临界常$5\times10^5$；圆管约2300/4000。别混。
\sect{“判断分离”}
看逆压梯度和壁面速度梯度。$\tau_w=\mu(\partial u/\partial y)_w$，分离点$\tau_w=0$。压差阻力因尾流低压而增大。

\chap{M1 往年题/课后题常见组合：压缩答案}
\sect{组合1：静水压+铰链}
题面有闸门、铰链、拉索：$F=\rho gh_CA$，$y_D=y_C+I_C/(y_CA)$，对铰链取矩：$T l_T=F l_F+G l_G$。若两侧有水，分别算两侧压力再相减。
\sect{组合2：曲面压力+浮力}
题面有圆弧闸门/半圆柱：水平分量=投影面压力；竖直分量=压力体重。若物体可浮起，加上浮力$\rho gV_{\rm disp}$和自重平衡。
\sect{组合3：U管+加速度}
题面容器加速或旋转：先写等效体力，压强梯度$\nabla p=\rho(\bm g-\bm a)$。自由液面垂直于等效重力。U管两端压差再用液柱关系。
\sect{组合4：速度场+势函数}
题面给$u,v$并问势/流：先算连续和旋度。若只满足连续，只有$\psi$；若只满足无旋，只有$\varphi$；两者都满足才都有。
\sect{组合5：速度场+应力}
先算速度梯度，再牛顿粘性应力：$\tau_{xy}=\mu(u_y+v_x)$，正应力含压强$-p+2\mu u_x$。若只给应力张量，直接$P\bm n$。
\sect{组合6：伯努利+动量}
喷嘴/弯管常先用伯努利求$V,Q$，再用动量求力。不要只写一个方程；速度未知时动量方程不能闭合。
\sect{组合7：泵+管路损失}
先由$Q,d$算速度，再$Re,\lambda$，再总损失，最后机械能方程。泵扬程和安装高度是不同问题：扬程给能量，安装高度受汽蚀限制。
\sect{组合8：水轮机+压力表}
压力表截面有$p/\rho g$和$V^2/(2g)$。若下游不是大水面，不能把速度头忽略。输出功率用水轮机取得的水头，不是总落差。
\sect{组合9：收缩喷管+背压}
如果$p_b/p_0$没有低到0.528，不能写阻塞。未阻塞出口压等于背压；阻塞出口压等于临界压，背压更低只影响管外调整。
\sect{组合10：Laval+面积比}
面积比给两个$M$，选择靠题意：喉前亚声速，喉后若设计超声速取超声速支；若全管亚声速取亚声速支。题中“无激波”很关键。
\sect{组合11：Laval+内激波}
有激波时不能从入口总压一路等熵到出口。必须波前$p_{01}$、过波总压损失、波后$p_{02}$。出口背压用于匹配激波位置。
\sect{组合12：斜激波+PM}
超声速绕多折角壁面：每个压缩角用斜激波，每个扩张角用PM。逐段更新$M,p,T$。最后若两支流相交，压强和方向需匹配。
\sect{组合13：边界层+平板阻力}
先判层/湍，再用平均摩擦系数求总阻力。若问局部壁剪，才用局部$c_f$。若问边界层厚度，取指定$x$处的$\delta(x)$。
\sect{组合14：湍流壁律+管内厚度}
给平均速度求$u_\tau$通常需要迭代；算出$u_\tau$后所有厚度都由$y^+=yu_\tau/\nu$反解。线性底层到5，缓冲到30。
\sect{组合15：Stokes小球+沉降}
小球终速题一定先写受力平衡：重力-浮力-阻力=0。若题给终速求黏度，反解$\mu$。最后核验$Re<1$。
\sect{组合16：相似准则}
模型试验要保持主要无量纲数相等。不可压重力主导：$Fr$相等；粘性主导：$Re$相等；高速气体：$Ma$相等。不能所有数都同时满足时说明近似。
\sect{组合17：空化}
局部压力最低处若$p_{\min}\le p_v$发生空化。伯努利题中速度最大点压强最低；泵吸入口高度越大，入口压越低。
\sect{组合18：流量计}
不可压流量计：压差$\to$速度$\to$流量；可压流量计：压差$\to M\to T,p,\rho,V\to \dot m$。可压缩不能直接用不可压Bernoulli。
\sect{组合19：水击/声速概念}
若题涉及扰动传播，液体声速$a=\sqrt{K/\rho}$，气体$a=\sqrt{\gamma RT}$。马赫数决定可压缩性重要程度，$M<0.3$常近似不可压。
\sect{组合20：自由涡}
自由涡$v_\theta=C/r$，环量$\Gamma=2\pi r v_\theta=2\pi C$常数。压强径向平衡$dp/dr=\rho v_\theta^2/r$，液面凹陷。
\sect{组合21：强迫涡}
强迫涡$v_\theta=\omega r$，像刚体旋转。$dp/dr=\rho\omega^2r$，自由面$z=\omega^2r^2/(2g)+C$。
\sect{组合22：管嘴出流}
大容器到喷口：$V=C_v\sqrt{2gH}$，$Q=C_dA\sqrt{2gH}$。若有管嘴损失，$H=V^2/(2g)+\sum h$。
\sect{组合23：多出口分流}
连续：$Q_{\rm in}=\sum Q_{\rm out}$。动量方程每个出口都要加$\rho Q_i\bm V_i$。若速度大小相同但方向不同，向量不能当标量相加。
\sect{组合24：压力恢复/扩压器}
亚声速扩压速度降、压力升；若逆压梯度过大，边界层分离，实际压力恢复低。超声速扩压需先激波减速。
\sect{组合25：雷诺平均}
速度$u=\bar u+u'$，$\overline{u'}=0$，但$\overline{u'v'}$不为0。RANS方程多出雷诺应力项，需要模型封闭。

\chap{M2 完整公式清单：按物理量查}
\sect{密度/压强/温度}
液体静压$p=p_a+\rho gh$。理想气体$p=\rho RT$。等温气体$\rho\propto p$。等熵气体$p/\rho^\gamma=C$，$T\rho^{1-\gamma}=C$。
\sect{速度/流量}
$Q=VA$，$\dot m=\rho Q$，圆管$A=\pi d^2/4$，环形$A=\pi(D^2-d^2)/4$。非均匀$Q=\int_A\bm V\cdot\bm n\,dA$。
\sect{水头}
总水头$H=p/\rho g+V^2/(2g)+z$。损失水头$h_L=h_f+\sum h_\zeta$。泵增$H_p$，水轮机取$H_T$。
\sect{沿程损失}
$h_f=\lambda(L/d)V^2/(2g)$。层流$\lambda=64/Re$。H--P给$\Delta p=32\mu LV/d^2$。湍流查Moody。
\sect{局部损失}
$h_\zeta=\zeta V^2/(2g)$。突扩损失$h=(V_1-V_2)^2/(2g)$。出口损失常$\zeta=1$，入口锐边约0.5，具体以题给为准。
\sect{动量通量}
一维均匀截面动量流率$\rho Q\bm V$。非均匀需$\int_A\rho\bm V(\bm V\cdot\bm n)dA$，常引入动量修正系数。
\sect{能量修正系数}
层流圆管动能修正系数$\alpha=2$；湍流近似$\alpha\approx1$。如果课程未强调，管路工程题通常取1，除非题给。
\sect{剪应力}
简单剪切$\tau=\mu du/dy$。圆管层流$\tau(r)=r\Delta p/(2L)$，壁面$\tau_w=R\Delta p/(2L)$。平板边界层$\tau_w=0.5\rho U^2 c_f$。
\sect{雷诺数}
圆管$Re=Vd/\nu$；平板$Re_x=Ux/\nu$；小球$Re=Ud/\nu$。特征长度随问题变，不要固定用直径。
\sect{马赫数和总量}
$M=V/\sqrt{\gamma RT}$，$T_0/T=1+(\gamma-1)M^2/2$，$p_0/p=(T_0/T)^{\gamma/(\gamma-1)}$，$\rho_0/\rho=(T_0/T)^{1/(\gamma-1)}$。
\sect{临界量}
$M=1$时$T^*/T_0=2/(\gamma+1)$，$p^*/p_0=[2/(\gamma+1)]^{\gamma/(\gamma-1)}$，空气为0.833和0.528。
\sect{面积--马赫}
$A/A^*=\frac1M[\frac{2}{\gamma+1}(1+\frac{\gamma-1}{2}M^2)]^{(\gamma+1)/(2\gamma-2)}$。注意指数为$(\gamma+1)/[2(\gamma-1)]$。
\sect{质量流量可压}
$\dot m=A p_0\sqrt{\gamma/(RT_0)}M(1+\frac{\gamma-1}{2}M^2)^{-(\gamma+1)/(2(\gamma-1))}$。空气阻塞系数$0.0404$。
\sect{正激波}
$M_2^2=\frac{1+(\gamma-1)M_1^2/2}{\gamma M_1^2-(\gamma-1)/2}$。$p_2/p_1=1+\frac{2\gamma}{\gamma+1}(M_1^2-1)$。$\rho_2/\rho_1=\frac{(\gamma+1)M_1^2}{(\gamma-1)M_1^2+2}$。
\sect{斜激波}
$M_{n1}=M_1\sin\beta$，法向按正激波，切向速度不变。$M_2=M_{n2}/\sin(\beta-\theta)$。$\theta-\beta-M$关系查表或公式。
\sect{PM函数}
$\nu(M)=\sqrt{\frac{\gamma+1}{\gamma-1}}\tan^{-1}\sqrt{\frac{\gamma-1}{\gamma+1}(M^2-1)}-\tan^{-1}\sqrt{M^2-1}$。膨胀角$\delta=\nu_2-\nu_1$。
\sect{平板边界层层流}
$\delta=5x/\sqrt{Re_x}$，$\delta^*=1.72x/\sqrt{Re_x}$，$\theta=0.664x/\sqrt{Re_x}$，$c_f=0.664/\sqrt{Re_x}$，$\bar C_f=1.328/\sqrt{Re_L}$。
\sect{平板边界层湍流}
$\delta=0.37x/Re_x^{1/5}$，$c_f=0.0592/Re_x^{1/5}$，$\bar C_f=0.074/Re_L^{1/5}$。混合转捩用减$1742/Re_L$修正。
\sect{边界层积分}
$\delta^*=\int_0^\infty(1-u/U)dy$，$\theta=\int_0^\infty(u/U)(1-u/U)dy$。无压梯度动量积分$\tau_w=\rho U^2d\theta/dx$。
\sect{Stokes}
$D=3\pi\mu dU$，$C_D=24/Re$。终速$U=(\rho_s-\rho)gd^2/(18\mu)$。
\sect{势流基本关系}
直角：$u=\phi_x=\psi_y,v=\phi_y=-\psi_x$。极坐标：$v_r=\phi_r=\psi_\theta/r$，$v_\theta=\phi_\theta/r=-\psi_r$。无旋不可压：$\nabla^2\phi=\nabla^2\psi=0$。
\sect{圆柱绕流}
$v_r=U(1-a^2/r^2)\cos\theta$，$v_\theta=-U(1+a^2/r^2)\sin\theta$。壁面$V=2U|\sin\theta|$，$C_p=1-4\sin^2\theta$。
\sect{环量升力}
$\Gamma=\oint\bm V\cdot d\bm l$。库塔--儒可夫斯基$L'=\rho U\Gamma$。点涡$v_\theta=\Gamma/(2\pi r)$。

\chap{M3 符号解释表：看到符号就知道干什么}
\sect{$\rho,\mu,\nu$}
$\rho$密度，$\mu$动力粘度，$\nu=\mu/\rho$运动粘度。Re题优先用$\nu$；剪应力题优先用$\mu$。单位：$\mu$为Pa·s，$\nu$为$m^2/s$。
\sect{$Q,\dot m,V,\bar u$}
$Q$体积流量，$\dot m=\rho Q$质量流量，$V$常为平均速度，$\bar u$明确表示截面平均/时均。圆管层流中心速度$u_{\max}=2\bar u$。
\sect{$p,p_g,p_0,p^*$}
$p$可能是绝压，$p_g$表压，$p_0$滞止压，$p^*$临界静压。气体状态方程必须绝压；水力方程可统一表压。
\sect{$A,A^*,A_t,A_e$}
$A$截面积，$A^*$临界面积，$A_t$喉部面积，$A_e$出口面积。阻塞时最小面积$A_t=A^*$；出口不一定是临界截面。
\sect{$h_f,h_\zeta,H_p,H_T$}
$h_f$沿程损失，$h_\zeta$局部损失，$H_p$泵扬程，$H_T$水轮机水头。能量方程中损失总在被消耗一侧。
\sect{$\lambda,\zeta,C_D,C_f$}
$\lambda$管摩阻系数，$\zeta$局部损失系数，$C_D$外绕流阻力系数，$C_f$摩擦阻力系数。四者不能混用。
\sect{$\delta,\delta^*,\theta$}
边界层名义厚度、位移厚度、动量厚度。$\delta$看速度达到$0.99U$，后两者由积分定义。
\sect{$u_\tau,y^+,u^+$}
摩擦速度$u_\tau=\sqrt{\tau_w/\rho}$；无量纲壁距$y^+=yu_\tau/\nu$；无量纲速度$u^+=u/u_\tau$。用于湍流近壁，不用于普通层流H--P。
\sect{$M,\mu,\beta,\theta$}
压缩流里$M$马赫数，$\mu$常作马赫角$\sin\mu=1/M$，$\beta$斜激波波角，$\theta$气流偏转角。不要把粘度$\mu$和马赫角$\mu$混。
\sect{$\nu(M)$}
Prandtl--Meyer函数，不是运动粘度$\nu$。膨胀波用$\nu(M)$，粘性Re用运动粘度$\nu$，看上下文。

\chap{M4 分支流程：从条件到公式}
\sect{流程A：给速度场}
1 算散度。2 算旋度。3 若散度0求$\psi$；若旋度0求$\phi$；若问流线用$dy/dx=v/u$；若问加速度用当地+迁移。4 若给粘度求应力，再速度梯度。
\sect{流程B：给静止液体图}
1 选自由液面为$p_a$。2 沿竖直高度算压强。3 平面用形心，曲面分水平/竖直。4 求铰链/拉力时对铰点取矩。
\sect{流程C：给管路图}
1 标点1、2，通常选大水面/压力表/出口。2 算各管段$V,Re,\lambda$。3 写能量方程含$h_f,h_\zeta,H_p/H_T$。4 解压力、扬程或功率。
\sect{流程D：给喷管图}
1 判断入口速度是否可忽略，定$p_0,T_0$。2 看是收缩还是Laval。3 用背压/面积比判断$M$。4 等熵段算状态；激波处分段。
\sect{流程E：给激波图}
1 看压缩角还是膨胀角。2 压缩角：斜激波；膨胀角：PM。3 斜激波先求$M_{n1}$。4 多波逐段更新，不能跨波共用总压。
\sect{流程F：给边界层/平板}
1 算$Re_x$或$Re_L$。2 判断层流/湍流/转捩。3 问厚度用$\delta$，问阻力用$\bar C_f$，问壁剪用局部$c_f$。4 面积看单面/双面。
\sect{流程G：给层流解析}
1 写N--S假设：定常、不可压、单向、充分发展。2 化成$\mu u''=$常数。3 积分两次。4 代无滑移/交界条件。5 求$Q,\tau,\Delta p$。
\sect{流程H：给相似/量纲}
1 列变量。2 写基本量纲。3 选重复变量。4 构造$\pi$群。5 说明控制参数，如$Re,Fr,Ma$。

\chap{M5 易错点逐条排雷}
\sect{错1：把背压当出口压}
收缩管未阻塞可令$p_e=p_b$；阻塞后$p_e=p^*$。Laval非设计状态出口外有波系，$p_e$和$p_b$也未必相等。
\sect{错2：把0.528到处用}
0.528是空气等熵$M=1$时$p^*/p_0$。只有截面临界才用。文丘里喉部若$M<1$不能用。
\sect{错3：激波后继续用原$p_0$}
错。正激波总温不变，总压下降。波后所有等熵计算都用$p_{02}$，不是$p_{01}$。
\sect{错4：面积比支路选错}
$A/A^*>1$有两个$M$。收缩管取亚声速；Laval扩张段根据背压和题意取亚/超声速。无激波设计通常取超声速。
\sect{错5：H--P用于湍流}
H--P抛物线只适用于层流充分发展圆管。湍流用摩阻系数/壁律，不能说$u_{\max}=2\bar u$。
\sect{错6：平板阻力用截面积}
外流平板摩擦阻力面积是湿表面积$Lb$或$2Lb$，不是迎风面积，也不是管截面积。
\sect{错7：压力中心公式乱用}
$y_D=y_C+I_C/(y_CA)$中的$y$是沿平面从自由液面交线量的坐标；竖直深度要乘$\sin\theta$。
\sect{错8：动量方程标量化}
弯管、斜板必须用向量分量。速度大小不变不代表动量不变，方向变也有力。
\sect{错9：局部损失速度取错}
局部元件在哪个管段，就用哪个管段速度。突扩/突缩常需要前后速度，不要全用主干管速度。
\sect{错10：表压绝压混用}
压力表读数是表压；汽化压和气体方程用绝压。能量方程若两点同大气，可用表压统一。
\sect{错11：PM压缩也用}
PM函数描述超声速等熵膨胀。超声速压缩通常是斜激波，不是把$\nu$相减就完。
\sect{错12：边界层厚度和管半径混}
管内湍流的$y$是离壁距离，$y=R-r$。边界层厚度不是管半径，近壁层只是一小部分。
\sect{错13：把流函数当势函数}
流线是$\psi=C$，等势线是$\phi=C$，二者正交。速度关系符号不同，特别是$v=-\psi_x$。
\sect{错14：源涡原点直接用}
点源、点涡在奇点处模型失效，速度/压强发散。题通常只问奇点外区域。
\sect{错15：忘记单位换算}
cm$^2$、mm、kPa、mm$^2$/s最容易错。写表前先统一SI单位。

\chap{M6 详细小模板：把答案写得像标准解}
\sect{标准解模板：静水压}
解：取自由液面为基准，形心深度$h_C=\cdots$。合力$F=\rho gh_CA=\cdots$。压力中心$y_D=y_C+I_C/(y_CA)=\cdots$。由力矩平衡$\sum M_O=0$得所求力$\cdots$。
\sect{标准解模板：动量}
解：取包围弯管/喷嘴的控制体。连续方程得$Q=AV$。动量方程$\sum F_x=\rho Q(V_{2x}-V_{1x})$，$\sum F_y=\rho Q(V_{2y}-V_{1y})$。解得壁面对流体力，所求管受力为反向。
\sect{标准解模板：管路}
解：选1、2截面列机械能方程。各段速度$V_i=Q/A_i$，$Re_i=V_id_i/\nu$，查得$\lambda_i$。总损失$\sum\lambda_iL_i/d_i\,V_i^2/(2g)+\sum\zeta_jV_j^2/(2g)$。代入得$H_p/H_T/p$。
\sect{标准解模板：N--S层流}
解：定常、不可压、充分发展、单向流，$u=u(y)$或$u(r)$。N--S化为$\mu u''=\cdots$。积分两次得通解。边界无滑移/剪应力连续求常数。最后求$Q,\tau$。
\sect{标准解模板：势流}
解：该流动不可压无旋，故存在$\phi,\psi$且满足Laplace方程。由基本流叠加得总$\psi$。固壁为流线$\psi=C$。由速度关系求$V$，由伯努利求$p$。
\sect{标准解模板：喷管}
解：入口/容器状态为滞止状态$p_0,T_0$。先判断$p_b/p_0$与临界压比或用面积比求$M$。等熵段用$p_0/p(M),T_0/T(M)$。若阻塞，$M=1$且$\dot m=\dot m_{\max}$。
\sect{标准解模板：激波}
解：波前由等熵或面积比求$M_1,p_1,T_1$。正激波/斜激波关系得$M_2,p_2,T_2$。总压损失后得$p_{02}$。波后再按等熵关系继续到出口。
\sect{标准解模板：边界层}
解：先算$Re_L$判别流态。按层流/湍流选$\delta(x),C_f$公式。局部量用$x$处局部公式，整体阻力用平均摩擦系数。代入湿表面积得$D$。
\sect{标准解模板：相似}
解：选取影响物理量$\{...\}$，基本量纲为$M,L,T$，故无量纲数个数$n-3$。选重复变量$\rho,V,L$，得$Re,Fr,\cdots$。模型与原型满足主导无量纲数相等。

\chap{M7 补充题型：自由面与涡}
\sect{自由涡液面}
$v_\theta=C/r$，径向平衡$dp/dr=\rho C^2/r^3$。自由面$p=p_a$且竖向静压，得$z=C_0-C^2/(2gr^2)$，中心凹陷。适用于排水旋涡外部近似。
\sect{强迫涡液面}
$v_\theta=\omega r$，$dp/dr=\rho\omega^2r$，自由面$z=\omega^2r^2/(2g)+C$，抛物面。容器旋转题用这个。
\sect{复合涡}
核心强迫涡、外部自由涡，交界处速度连续。常用于实际旋涡近似。压强分布分段积分。
\sect{Rankine卵形}
均匀流+源汇对。物体边界为过驻点的闭合流线。长宽由驻点位置和分界流线确定，考试通常只需写总$\psi$并说明。
\sect{偶极子物理意义}
源汇距离趋零、强度趋无穷但乘积有限，得到偶极。圆柱绕流就是均匀流+偶极，偶极强度$M=Ua^2$。
\sect{机翼升力}
势流中机翼可由保角变换把圆柱映到翼型；环量由Kutta条件确定；升力公式仍为$L'=\rho U\Gamma$。考试概念题写这个即可。

\chap{M8 最后一页应该带的索引}
\sect{教材查找建议}
正激波表、斜激波表、PM表、面积--马赫表、Moody图、阻力系数图、常见截面惯性矩表、物性表。速查表里不要抄全表，只写“什么时候查、查完怎么用”。
\sect{考试先后顺序}
先做能套模板的：静水压、速度场、管路、喷管、边界层。再做需查表的激波/PM。最后做推导证明题。不会的题先写模型和控制方程拿步骤分。
\sect{数值题通用检查}
压力不能低于绝对零；水泵入口绝压不能随便小于汽化压；质量流量在同一管内守恒；激波后静压上升；膨胀后静压下降；边界层厚度不应超过特征长度太多。
\sect{概念题通用答法}
先定义，再说物理原因，再说公式或后果。例如问阻塞：定义$M=1$，原因信息不能上游传递，后果质量流量最大且不随背压增大。
\sect{证明题通用答法}
从守恒律或定义出发。连续方程来自质量守恒；Euler来自无粘动量；Bernoulli来自沿流线积分Euler；势流Laplace来自无旋+不可压；流函数Poisson来自涡量定义。
\sect{图像题通用答法}
先在图上标1、2截面、坐标方向、速度方向、压力方向、已知高度。图像题扣分多来自方向和截面选错。
\sect{答题纸三行保底}
1 “取控制体/流线/截面如图。” 2 “由连续/动量/能量/状态方程得...” 3 “代入题给数据，取正方向，故结果为...” 这三行至少保住结构分。

\chap{M9 最后一页：高频小问速答}
\sect{为什么应力张量对称}
对微元取角动量矩平衡，若无体偶力，剪应力互等，故应力张量对称。用途：减少未知量，保证任意切面的应力由$P\bm n$唯一确定。
\sect{为什么有流函数}
二维不可压连续方程$u_x+v_y=0$保证存在$\psi$使$u=\psi_y,v=-\psi_x$。流函数差等于两条流线间单位宽流量。
\sect{为什么有势函数}
无旋$\nabla\times\bm V=0$且区域单连通时，速度场为某标量势的梯度$\bm V=\nabla\phi$。等势线与流线正交。
\sect{为什么Laplace}
不可压$\nabla\cdot\bm V=0$，又$\bm V=\nabla\phi$，所以$\nabla^2\phi=0$。无旋二维流函数同理$\nabla^2\psi=0$。
\sect{为什么Bernoulli沿流线}
无粘Euler方程沿流线积分，定常不可压时动能、压能、位能总和守恒。若无旋，常数可扩展到全场。
\sect{为什么圆柱理想阻力为0}
势流无粘且前后对称，压力积分$x$方向抵消。该结论与现实矛盾，称达朗贝尔佯谬，现实阻力由粘性分离导致。
\sect{为什么边界层存在}
高Re外流粘性只在近壁薄层内重要；壁面无滑移使速度从0快速过渡到外流速度，形成边界层。
\sect{为什么逆压梯度会分离}
流向压力升高使近壁低能流体减速，壁面速度梯度降至0并反向，流线脱离壁面形成回流区。
\sect{为什么湍流阻力大}
脉动增强动量交换，等效剪应力增大；边界层更厚但抗分离能力更强。工程上用经验摩阻系数。
\sect{为什么阻塞}
喉部$M=1$后，下游压力扰动不能逆流传到上游；在给定$p_0,T_0,A^*$下质量流量达到最大。
\sect{为什么激波总压下降}
激波虽绝热，但不可逆熵增。滞止温度由能量守恒不变，滞止压力因熵增下降。
\sect{为什么膨胀波等熵}
膨胀扇由无数弱扰动连续组成，近似无耗散、无间断，故可按等熵关系处理。
\sect{为什么斜激波用法向马赫数}
激波面切向无压力梯度，切向速度不变；压缩只发生在法向方向，所以法向分量满足正激波关系。
\sect{为什么Laval能超声速}
亚声速在收缩段加速至喉部$M=1$；过喉后若背压合适，超声速在扩张段继续加速。单纯收缩管不能稳定得到超声速出口。
\sect{为什么粗糙度层流不用}
层流摩阻来自粘性剪切解析解，$\lambda=64/Re$；粗糙凸起若未破坏层流，公式中不出现粗糙度。湍流中粗糙度影响近壁结构。
\sect{为什么平板局部/平均系数不同}
局部$c_f(x)$随$x$变化；总阻力是沿全板积分的平均效果，所以用$\bar C_f$。尾部局部值不能代表总板阻力。
\sect{为什么两层流体剪应力连续}
交界面没有质量和表面力奇异时，切向力平衡要求两侧剪应力相等；同时无滑移要求速度连续。
\sect{为什么H--P抛物线}
圆管层流充分发展时压力梯度为常数，粘性扩散方程在圆柱坐标积分两次，边界无滑移，得到$r^2$抛物线。
\sect{为什么Stokes阻力与速度一次方成正比}
低Re时惯性项可忽略，N--S线性化，粘性力主导，因此阻力与速度成正比而不是平方。
\sect{为什么管路损失用速度头}
损失本质是机械能耗散，工程上以单位重量能量表示。局部扰动强度通常与动压$\rho V^2/2$成正比，所以写成$\zeta V^2/(2g)$。
\sect{为什么压力中心低于形心}
静水压随深度线性增大，较深部分受力更大，合力作用点向深处偏移，故压力中心低于形心。
\sect{为什么动量方程要控制体}
流体不断进出，质点系统不好追踪。控制体形式用流入流出动量通量表达，适合弯管、喷嘴、射流等定常问题。
\sect{为什么流量要积分}
截面速度若不均匀，$VA$只适用于平均速度。真实流量是每个面积元的法向速度贡献积分：$Q=\int_A v_n dA$。
\sect{为什么表压可为负}
表压相对大气压定义，低于大气压时为负；绝压不能为负。汽蚀判断看绝压是否低于汽化压。
\sect{为什么马赫数重要}
马赫数表示速度与声速之比，决定压缩性和扰动传播范围。$M<1$扰动可上游传播，$M>1$只能在马赫锥内传播。
\sect{为什么扩压器易分离}
扩压器亚声速减速升压，存在逆压梯度；若扩张角过大，边界层不能克服逆压梯度，产生分离，压力恢复下降。
\sect{为什么泵入口最危险}
泵吸入侧压力最低，若安装高度大或损失大，入口绝压可能低于汽化压，先发生汽蚀。
\sect{为什么水轮机要扣损失}
上游总水头不是全部转化为轴功，一部分作为沿程/局部损失和下游剩余动能，剩余才是水轮机可利用水头。
\sect{为什么圆管湍流用壁律}
湍流N--S时均后出现未知雷诺应力，不能直接解析闭合；近壁实验显示无量纲速度分布具有通用壁律。
\sect{为什么位移厚度叫位移}
边界层内速度亏损使外流好像被壁面向外排挤一段距离，等效距离就是$\delta^*$。
\sect{为什么动量厚度重要}
动量厚度表示边界层造成的动量亏损，动量积分方程把壁剪和$\theta$的增长联系起来。
\sect{为什么答案要写假设}
流体题公式都有适用条件。写“定常、不可压、无粘/充分发展/等熵”等假设，可以明确选用方程并获得过程分。

\chap{M10 速算区：不用重新推的中间结果}
\sect{常用平方/根号}
$\sqrt{2gH}\approx4.43\sqrt{H}$ m/s。$V^2/(2g)$：$V=1,2,5,10$ m/s时约$0.051,0.204,1.27,5.10$m。
\sect{常用水柱压强}
$1$m水柱约$9.81$kPa；$10$m约$98.1$kPa；$10.33$m约1atm。水银$1$mm约133Pa。
\sect{常用圆管面积}
$d=0.1$m，$A=7.85\times10^{-3}$；$d=0.2$m，$A=3.14\times10^{-2}$；$d=0.3$m，$A=7.07\times10^{-2}$。
\sect{空气等熵速算}
$M=0.5$：$T_0/T=1.05,p_0/p\approx1.19$。$M=1$：$1.2,1.893$。$M=2$：$1.8,7.82$。$M=3$：$2.8,36.7$。
\sect{空气临界速算}
$p^*/p_0=0.528$，$T^*/T_0=0.833$，$\rho^*/\rho_0=0.634$。$a^*=\sqrt{\gamma RT^*}$，$V^*=a^*$。
\sect{正激波空气$M_1=2$}
$M_2\approx0.577$，$p_2/p_1=4.5$，$\rho_2/\rho_1=2.67$，$T_2/T_1=1.69$。可用于检查数量级。
\sect{正激波空气$M_1=3$}
$M_2\approx0.475$，$p_2/p_1\approx10.33$，$\rho_2/\rho_1\approx3.86$，$T_2/T_1\approx2.68$。
\sect{层流圆管速算}
$Re=VD/\nu$。水$\nu\approx10^{-6}$，所以$V=1,d=0.1$时$Re=10^5$湍流；油$\nu=10^{-5}$时同条件$Re=10^4$。
\sect{H--P压降速算}
$\Delta p=32\mu LV/d^2$。直径减半，压降增为4倍（同平均速度）；同流量下$V\sim1/d^2$，压降$\sim1/d^4$。
\sect{平板层流厚度速算}
$\delta/x=5/\sqrt{Re_x}$。$Re_x=10^5$，$\delta/x=0.0158$；$Re_x=10^6$，$\delta/x=0.005$。
\sect{平板层流阻力速算}
$\bar C_f=1.328/\sqrt{Re_L}$。$Re_L=10^5$时$0.0042$；$10^6$时$0.00133$。
\sect{湍流平板阻力速算}
$\bar C_f=0.074/Re_L^{1/5}$。$Re_L=10^6$时约$0.0047$；$10^7$时约$0.00295$。
\sect{壁律厚度速算}
若$u_\tau=0.1$m/s，水$\nu=10^{-6}$，线性底层$y_5=5\times10^{-5}$m，缓冲层外$y_{30}=3\times10^{-4}$m。
\sect{Stokes终速速算}
$U=(\rho_s-\rho)gd^2/(18\mu)$。直径增10倍，终速增100倍；但需重新检查$Re<1$。
\sect{动量力速算}
水流$Q=0.01$m$^3$/s，$V=10$m/s，动量力$\rho QV=100$N。速度方向转90度，两个方向分量都要算。
\sect{泵功率速算}
$P=\rho gQH$。水中$Q=0.01,H=10$m，水功率约981W。考虑效率0.7，轴功率约1.4kW。
\sect{粗糙度速算}
$\varepsilon=0.2$mm，$d=0.2$m，则$\varepsilon/d=0.001$。题中“粗糙度$h$”通常就是$\varepsilon$。
\sect{声速速算}
空气$T=293$K，$a=\sqrt{1.4\times287\times293}\approx343$m/s。温度越高声速越大，$a\propto\sqrt T$。
\sect{马赫角速算}
$M=2$，$\mu=30^\circ$；$M=3$，$\mu\approx19.5^\circ$；$M=1.5$，$\mu\approx41.8^\circ$。
\sect{PM函数使用提醒}
PM函数角度从$M=1$时0开始，随$M$增大而增大。空气$M=2$时$\nu\approx26.4^\circ$，$M=3$时$\nu\approx49.8^\circ$，可查数量级。
\sect{斜激波弱解检查}
弱斜激波后通常仍可能超声速；强解后多为亚声速。工程外流默认弱解。若算出波角接近$90^\circ$，接近正激波。
\sect{喷管状态数量级}
出口超声速时静压通常远低于总压；若算得$p_e>p_0$必错。等熵加速时$T,p,\rho$都下降，$V,M$上升。
\sect{管路答案数量级}
损失水头不应为负；泵扬程若小于总损失无法供水；水轮机可利用水头不能超过上下游总水头差。
\sect{压力中心数量级}
矩形竖直板上边在自由面、高$h$，压力中心深度$2h/3$。若算到$h/2$上方，说明公式或坐标错。
\sect{圆管层流剪应力}
壁面剪应力$\tau_w=8\mu\bar u/d$。直径越小剪应力越大；粘度越大剪应力越大。
\sect{局部损失数量级}
$\zeta=1,V=2$m/s时损失约$0.204$m水头；$V=10$m/s时约$5.1$m。高速下局部损失很大。
\sect{不可压近似}
气体$M<0.3$时密度变化通常小于约5\%，可近似不可压；液体通常不可压，除非题明确弹性模量/水击。
\sect{题目读数顺序}
先圈单位，再圈“忽略/不可忽略”，再圈“等熵/有激波/有损失”，最后圈要求量。很多错误来自漏掉“入口速度可忽略”或“表压”。
\sect{考前最后记忆}
静水压看高度；动量看方向；管路看损失；势流看叠加；喷管看$M$和背压；激波看分段；边界层看$Re$；粘性解析看边界条件。

\chap{M11 展开型标准答案：可直接仿写}
\sect{例：速度场完整答}
给$\bm V=(u,v,w)$。先写$\nabla\cdot\bm V=u_x+v_y+w_z=\cdots$，故是否不可压。再写$\nabla\times\bm V=(w_y-v_z,u_z-w_x,v_x-u_y)=\cdots$，故是否有旋。若无旋，$\phi=\int u dx+f(y,z)$，逐项定函数。
\sect{例：流线完整答}
由流线微分方程$\frac{dx}{u}=\frac{dy}{v}$，代入$t=t_0$若非定常。整理为$M(x,y)dx+N(x,y)dy=0$，积分得$F(x,y)=C$。若与$\psi=C$一致，说明计算正确。
\sect{例：迹线完整答}
由$\dot x=u(x,y,t),\dot y=v(x,y,t)$。先解$x(t)$，再代入解$y(t)$；带入初始条件确定常数。最后消去$t$或写参数方程。
\sect{例：应力完整答}
单位法向先归一化。$\bm t=P\bm n$写成列向量。$\sigma_n=\bm n^T P\bm n$。$\tau=\sqrt{\bm t^T\bm t-\sigma_n^2}$。若题要分量，写$\bm t_N=\sigma_n\bm n,\bm t_T=\bm t-\bm t_N$。
\sect{例：闸门完整答}
设水深方向坐标$y$，形心$y_C$，面积$A$，惯性矩$I_C$。$F=\rho g y_C\sin\theta A$。压力中心$y_D=y_C+I_C/(y_CA)$。对铰点取矩求未知力，注意力臂是几何垂距。
\sect{例：曲面完整答}
投影面积$A_x$形心深$h_C$，$F_x=\rho gh_CA_x$。压力体体积$V_p$，$F_y=\rho gV_p$，方向由水压推曲面方向定。合力大小和方向$\tan\alpha=F_y/F_x$。
\sect{例：弯管完整答}
先用连续得$Q,A,V$。假设支座对流体反力$R_x,R_y$。列$x,y$动量，包含压力力$pA$和重力（若不可忽略）。解得$R$，题问管道受力则$-R$。
\sect{例：泵完整答}
从吸水池水面到泵入口：求安装高度/汽蚀。从泵入口到出口或两水池：求泵扬程。不要把两个方程混成一个。汽蚀题必须用绝压。
\sect{例：水轮机完整答}
选上游高水面和下游压力表/水面。机械能方程右侧加$H_T$，即流体失去的水头。$H_T$求出后$P=\eta\rho gQH_T$。
\sect{例：圆管层流完整答}
由受力或N--S得$\frac1r\frac{d}{dr}(rdu/dr)=\frac1\mu dp/dx$。有限性使$r=0$处常数不发散，壁面$u(R)=0$。得到$u(r)$后积分求$Q$。
\sect{例：平板Couette完整答}
方程$\mu u''=dp/dx$。通解$u=(dp/dx)y^2/(2\mu)+C_1y+C_2$。边界$u(0)=0,u(h)=U$。若给流量或中点速度，再解$dp/dx$或$U$。
\sect{例：两层流完整答}
两层各有$u_i=A_iy+B_i$。四条件：下壁、上壁、交界速度连续、交界剪应力连续。若无压差，剪应力常数，直接用阻力串联公式。
\sect{例：势流叠加完整答}
列出各基元$\psi_i$，总$\psi=\sum\psi_i$。物体边界通常为某条$\psi=C$。驻点令$\nabla\phi=0$或速度分量为0。压力由伯努利，合力由积分。
\sect{例：圆柱有环量完整答}
总速度壁面$v_\theta=-2U\sin\theta+\Gamma/(2\pi a)$。停滞点令$v_\theta=0$。压力系数$C_p=1-(v_\theta/U)^2$。升力用$L'=\rho U\Gamma$。
\sect{例：Laval喷管完整答}
写明“喉部达到临界则$A_t=A^*$”。面积比求出口$M_e$，等熵求$p_e$。比较$p_e$与$p_b$判断设计/过膨胀/欠膨胀。若管内激波，分三段。
\sect{例：正激波完整答}
波前状态已知或由等熵求。用正激波表得$M_2,p_2/p_1,T_2/T_1$。若要波后速度，$a_2=\sqrt{\gamma RT_2}$，$V_2=M_2a_2$。
\sect{例：斜激波完整答}
已知$M_1,\theta$查$\beta$。$M_{n1}=M_1\sin\beta$，正激波关系得$M_{n2}$。$M_2=M_{n2}/\sin(\beta-\theta)$。压力温度比与法向马赫数相关。
\sect{例：膨胀波完整答}
查$\nu(M_1)$，加转角得$\nu(M_2)$，反查$M_2$。再等熵求$p_2,T_2$。若求壁面角，膨胀扇总角等于气流偏转角。
\sect{例：平板边界层完整答}
$Re_L=UL/\nu$。判层流/湍流。阻力$D=\bar C_f(0.5\rho U^2)A$，尾部厚度$\delta(L)$。若双面流动，面积乘2。
\sect{例：动量积分边界层}
假设$u/U=f(y/\delta)$。算$\theta=C\delta$和$\tau_w=\mu U f'(0)/\delta$。代$\tau_w=\rho U^2 d\theta/dx$，分离变量积分得$\delta(x)$。
\sect{例：Stokes小球完整答}
写$Re=\rho Ud/\mu<1$，故用Stokes。受力平衡：$G-F_B-D=0$。代$G=\rho_sg\pi d^3/6$，$F_B=\rho g\pi d^3/6$，$D=3\pi\mu dU$。

\chap{M12 最终保底：不会算也能写的东西}
\sect{压强题保底}
写“压强随竖直深度线性变化，与容器形状无关”。再画压力分布三角形/梯形。即使不会算力矩，也能拿基本分。
\sect{速度场题保底}
写散度、旋度两个算式。多数题至少一半分在这里。若能积分出$\phi$或$\psi$，再补。
\sect{管路题保底}
写完整机械能方程，哪怕$\lambda$不会查。把$h_f,\sum h_\zeta,H_p,H_T$位置放对，能拿结构分。
\sect{喷管题保底}
写等熵三公式和临界0.528；再写“先判阻塞，再算$M$”。不要直接乱代数。
\sect{激波题保底}
写“波前等熵、过波用激波关系、波后用新总压等熵”。这句比一堆错公式值钱。
\sect{边界层题保底}
写$Re_x,Re_L$，选层流/湍流，写$\delta,C_f,D$三类公式。题问什么就取对应公式。
\sect{粘性解析题保底}
写N--S简化假设和边界条件。即使积分错，假设和边界条件也有分。
\sect{概念题保底}
定义+原因+后果三段式。例：阻塞=喉部$M=1$；原因=扰动不能上传；后果=质量流量最大。
\sect{证明题保底}
从定义写起：连续来自质量守恒，伯努利来自Euler沿流线积分，Laplace来自无旋不可压，流函数Poisson来自涡量定义。
\sect{单位题保底}
所有数值先写SI单位转换。老师看见单位转换，至少知道你没瞎代。

\chap{M13 专题微题：一题一行解法}
\sect{微题1：水面下矩形窗}
$F=\rho g(h_0+h/2)bh$；压力中心距窗顶$y=(h_0+h/2)+I_C/[(h_0+h/2)A]$，再换成题目坐标。
\sect{微题2：圆形舱门}
$A=\pi R^2,I_C=\pi R^4/4$。$F=\rho gh_CA$，$y_D=y_C+R^2/(4y_C)$。
\sect{微题3：浮筒吃水}
平衡$\rho_w gV_{\rm disp}=G$。若圆柱竖放，$V_{\rm disp}=Ah$；若横放，查/算圆缺面积。
\sect{微题4：水箱小孔时间}
$A_t dH/dt=-C_dA_o\sqrt{2gH}$，积分$t=\frac{2A_t}{C_dA_o\sqrt{2g}}(\sqrt{H_1}-\sqrt{H_2})$。
\sect{微题5：皮托静压管}
$p_0-p=\frac12\rho V^2$。若接U管，$p_0-p=(\rho_m-\rho)g\Delta h$。
\sect{微题6：文丘里不可压}
$Q=C_dA_2\sqrt{2\Delta p/[\rho(1-\beta^4)]}$，$\beta=d_2/d_1$。
\sect{微题7：喷嘴反力}
$F_{\rm hold}\approx\rho QV+(p_e-p_a)A_e$，若出口对大气$p_e=p_a$，只剩动量项。
\sect{微题8：平板拖曳}
层流单面$D=1.328/\sqrt{Re_L}\cdot0.5\rho U^2Lb$。双面乘2。
\sect{微题9：球低Re阻力}
$D=3\pi\mu dU$；若给$C_D$，低Re可写$24/Re$。
\sect{微题10：圆管层流流量}
$Q=\pi R^4\Delta p/(8\mu L)$。半径误差影响四次方，注意用$R$不是$d$。
\sect{微题11：管内平均速度}
$\bar u=Q/A$。层流最大速度$2\bar u$；湍流最大速度与平均速度关系需经验式，不用2。
\sect{微题12：局部突扩}
$h=(V_1-V_2)^2/(2g)$。压力可能升高但总能量下降。
\sect{微题13：泵效率}
$\eta=\rho gQH/P_{\rm shaft}$。若求轴功率，$P_{\rm shaft}=\rho gQH/\eta$。
\sect{微题14：水轮机效率}
$P_{\rm out}=\eta\rho gQH_T$。若给输出功率反求$H_T=P/(\eta\rho gQ)$。
\sect{微题15：自由涡压差}
$dp/dr=\rho C^2/r^3$，积分$p_2-p_1=\frac12\rho C^2(1/r_1^2-1/r_2^2)$。
\sect{微题16：强迫涡压差}
$dp/dr=\rho\omega^2r$，积分$p_2-p_1=\frac12\rho\omega^2(r_2^2-r_1^2)$。
\sect{微题17：声速}
空气$a\approx20.05\sqrt{T}$，$T$用K。$T=288$K时$a\approx340$m/s。
\sect{微题18：滞止温度}
$T_0=T+V^2/(2c_p)$。理想气$c_p=\gamma R/(\gamma-1)$。与$T_0/T=1+0.2M^2$等价。
\sect{微题19：收缩管最大流量}
空气$\dot m_{\max}=0.0404A p_0/\sqrt{T_0}$，$p_0$Pa，$A$m$^2$，$T_0$K，结果kg/s。
\sect{微题20：激波后速度}
先求$T_2,M_2$，再$V_2=M_2\sqrt{\gamma RT_2}$。不要用$V_1$直接乘压力比。
\sect{微题21：PM膨胀压力比}
$p_2/p_1=[(1+0.2M_1^2)/(1+0.2M_2^2)]^{3.5}$。
\sect{微题22：斜激波波角}
波角$\beta$大于马赫角$\mu$，小于$90^\circ$。若算出$\beta<\mu$，必错。
\sect{微题23：边界层位移厚度}
若$u/U=y/\delta$，$\delta^*=\delta/2$。若$u/U=(y/\delta)^{1/7}$，$\delta^*=\delta/8$。
\sect{微题24：边界层动量厚度}
线性剖面$\theta=\delta/6$；抛物$2\eta-\eta^2$时$\theta=2\delta/15$。
\sect{微题25：壁面剪应力}
由速度剖面求$\tau_w=\mu(\partial u/\partial y)_{y=0}$。边界层题不要用管流$\lambda$，除非题是管内。
\sect{微题26：两板间流量}
$q=\int_0^hu\,dy$，若$u=Uy/h$，$q=Uh/2$每单位宽。
\sect{微题27：两层流体剪力}
等效“粘性电阻”$h_1/\mu_1+h_2/\mu_2$，剪应力$\tau=U/$总阻力。
\sect{微题28：雷诺应力符号}
RANS中$-\rho\overline{u'v'}$像附加剪应力。若$\overline{u'v'}<0$，雷诺剪应力为正。
\sect{微题29：模型律}
若重力波模型保持$Fr$，速度比$V_m/V_p=\sqrt{L_m/L_p}$。时间比$t_m/t_p=\sqrt{L_m/L_p}$。
\sect{微题30：终极检查}
每个答案后写单位；每个力题说明方向；每个压缩流题说明是等熵还是激波；每个粘性题说明层流还是湍流。

\chap{M14 第六页补全：专题判断与查表动作}
\sect{喷管查表动作1}
已知$p_0/p$求$M$：查等熵表或解$(1+0.2M^2)^{3.5}=p_0/p$。若解得$M<1$但题说超声速，说明给的不是该段等熵关系或有激波。
\sect{喷管查表动作2}
已知$A/A^*$求$M$：表中读两支。先在题边写“亚支/超支”，避免抄错。出口扩张段常需要超支，入口收缩段常亚支。
\sect{激波查表动作}
正激波表读数后，立刻写$p_{02}=p_{01}(p_{02}/p_{01})$，后续所有出口压力用$p_{02}$。这是最容易漏的步骤。
\sect{斜激波查表动作}
$M_1,\theta\to\beta$。若表无精确值，用线性插值。读到$\beta$后不要忘了$M_{n1}=M_1\sin\beta$。
\sect{PM查表动作}
先查$\nu_1$，加膨胀角得$\nu_2$，再反查$M_2$。若题给压强比，也可先等熵求$M_2$，再求角度。
\sect{Moody查图动作}
先判层流/湍流。湍流才用$Re$和$\varepsilon/d$查图。若题给“完全粗糙”，$\lambda$几乎与$Re$无关。
\sect{阻力图查图动作}
先确定特征面积和特征长度。球/圆柱$A$为迎风面积；平板摩擦用湿表面积。查到$C_D$后$D=C_DqA$。
\sect{物性表查表动作}
水温给出时查$\rho,\mu,\nu,p_v$。汽蚀题必须查$p_v$。空气题通常给$R,\gamma$，若不给取空气标准值。
\sect{惯性矩查表动作}
压力中心题需要面积二次矩$I_C$，不是质量转动惯量。矩形$bh^3/12$，圆$\pi R^4/4$，三角形查表。
\sect{三角函数角度}
激波/PM表多为度；计算器确认DEG。LaTeX/推导里弧度不影响公式，但数值代入会错。

\chap{M15 题目图像再识别}
\miniimg{0.98}{assets/original_problem_schematics.png}
\sect{图像：有大水面+管道+泵}
一定是机械能方程。大水面速度0、表压0。泵前若问安装高度，就是汽蚀；泵后若问压力/功率，就是扬程。
\sect{图像：管道变径}
连续先给$V_1,V_2$。突扩有局部损失$(V_1-V_2)^2/(2g)$；渐扩看是否给扩压器效率或损失系数。
\sect{图像：管道弯头}
局部损失$\zeta V^2/(2g)$用于能量；弯管受力用动量。两个问题不要混。
\sect{图像：喷口射流}
自由射流出口压力等于大气。速度由伯努利；反力由动量。射流打板再按板方向分解动量。
\sect{图像：Laval喷管压力曲线}
压力曲线突然跳升表示正激波；平滑下降为等熵加速；出口外波系由$p_e$与$p_b$比较决定。
\sect{图像：楔形物体}
超声速来流撞楔尖，前方斜线是斜激波，楔角是流动偏转角$\theta$，斜线角是波角$\beta$。
\sect{图像：凹角/凸角}
超声速流遇凹角压缩，画斜激波；遇凸角膨胀，画膨胀扇。亚声速不画马赫波。
\sect{图像：平板外流}
前缘开始发展边界层。题问尾部厚度用$L$；题问阻力用全长平均摩擦系数；题问局部剪应力用$x$处局部系数。
\sect{图像：圆柱/球外绕}
低Re可能Stokes；高Re看分离和阻力系数。理想势流只给压力分布和升力概念，不给真实阻力。
\sect{图像：两平板夹流体}
上板动下板静是Couette；有压差是Poiseuille；两者同时存在则叠加。两层流体加交界条件。
\sect{图像：同心圆管/环隙}
若内外圆柱相对运动，是环隙Couette，速度分布可能含$\ln r$。若压差驱动，是环隙Poiseuille，查教材公式或由圆柱坐标积分。
\sect{图像：旋转容器}
自由面抛物线。若问压强，$p/\rho+gz-\omega^2r^2/2=C$。若问液面高度差，$\Delta z=\omega^2(R_2^2-R_1^2)/(2g)$。
\sect{图像：自由涡水面}
液面中心凹陷，$v_\theta=C/r$。与强迫涡相反：强迫涡中心低但外缘高为抛物面，自由涡近中心理论发散。
\sect{图像：U形管振荡}
两个自由面高度差产生回复力，周期$2\pi\sqrt{L/(2g)}$。若管截面不等，需要用体积连续换算位移。

\chap{M16 考试策略：三张纸怎么用}
\sect{第一页用途}
先看总判题、静水压、动量、能量。遇到任何题，先从第一页确定主方程。
\sect{第二页用途}
势流和理想流体。遇到$\phi,\psi$、圆柱、镜像、机翼、环量，直接翻这里。
\sect{第三页用途}
可压缩等熵和喷管。遇到$M,p_0,T_0,A^*$，先翻这里，不要先翻激波。
\sect{第四页用途}
激波、膨胀波和Laval复杂状态。确定有间断/楔角/膨胀角后翻这里。
\sect{第五页用途}
粘性管流、泵、水轮机、N--S解析。遇到$\mu,\nu,Re,\lambda$先翻这里。
\sect{第六页用途}
边界层、错题排雷、速算、保底模板。最后检查单位、方向、适用条件。
\sect{考场流程}
先用10秒归类，再用30秒写主方程，再代数值。卡住时不要停：写假设、写控制方程、写查表步骤。
\sect{拿分优先级}
过程分通常来自：模型判断、主方程、边界条件、查表/代入路径、单位方向。数值算错但路径对仍能得分；空白没有分。
\sect{不能承诺的事}
速查表不能替你读题和选择物理模型；但它能把模型到公式的路径固定下来。考试前至少拿一套题按本表跑一遍。

\chap{M17 章节最后核对清单}
\sect{第1章物性}
密度、粘度、压缩性、表面张力。会写$\nu=\mu/\rho$，$E=dp/(d\rho/\rho)$，$a=\sqrt{K/\rho}$或$\sqrt{\gamma RT}$。
\sect{第2章静力学}
静压强只随竖直高度变；平面压力用形心和压力中心；曲面压力分水平投影和压力体；相对平衡用等效重力。
\sect{第3章运动学}
会区分流线、迹线、脉线；会算加速度、散度、旋度；会判断有无流函数和势函数。
\sect{第4章动力学}
会用系统/控制体思想；会列连续、动量、能量；会算应力矢量$P\bm n$；会用Bernoulli及其限制。
\sect{第5章理想流体}
会写Euler、Bernoulli、势流基本方程；理解无粘不代表无压强；理解达朗贝尔佯谬。
\sect{第6章势流叠加}
会写均匀流、源、汇、涡、偶极；会叠加；会镜像；会圆柱绕流；会环量升力。
\sect{第7章可压缩}
会算马赫数、滞止量、等熵关系、临界状态、面积马赫、正/斜激波、PM膨胀、Laval背压状态。
\sect{第8章粘性流}
会N--S简化、H--P、Couette、两层流、管路损失、湍流壁律、边界层厚度、阻力和分离。
\sect{第9章湍流概念}
若考概念：时均+脉动分解，雷诺应力，DNS/LES/RANS，混合长度，对数律。
\sect{第10章边界层/绕流}
平板层流/湍流公式，位移/动量厚度，分离判断，外绕流阻力系数，Stokes小球。

\chap{M18 往年题可能问法：一句话答案}
\sect{问：流体和固体区别}
流体不能承受静止剪应力，任意小剪应力都会产生连续变形；固体可承受剪应力保持静止形变。
\sect{问：牛顿流体}
剪应力与速度梯度线性相关，$\tau=\mu du/dy$。水和空气常近似牛顿流体。
\sect{问：连续介质假设}
把分子离散结构在宏观上平均，物理量可看作连续函数。适用于尺度远大于分子平均自由程。
\sect{问：不可压条件}
严格为随体密度不变$D\rho/Dt=0$；常密度时$\nabla\cdot\bm V=0$。
\sect{问：流函数物理意义}
二维不可压中，两条流线间流函数差等于单位宽流量。
\sect{问：速度势物理意义}
势函数梯度等于速度；等势线与流线正交；存在条件是无旋。
\sect{问：环量}
速度沿闭合曲线的线积分$\Gamma=\oint\bm V\cdot d\bm l$；与升力和涡有关。
\sect{问：涡量}
涡量是速度旋度，表示微团旋转强度的两倍。无旋流可有环量奇点，如点涡外部。
\sect{问：Bernoulli常数}
沿同一流线守恒；若流动无旋，整个连通区域同一常数；实际流动需加损失。
\sect{问：动量修正系数}
非均匀速度截面动量通量不是$\rho Q\bar V$，需修正系数。层流圆管常较大，湍流接近1。
\sect{问：能量修正系数}
动能项用$\alpha V^2/(2g)$修正速度分布。层流圆管$\alpha=2$，湍流约1。
\sect{问：雷诺数意义}
惯性力与粘性力之比。小Re粘性主导，大Re惯性主导，决定层流/湍流和相似准则。
\sect{问：马赫数意义}
流速与声速之比。决定压缩性和扰动传播；超声速有马赫锥。
\sect{问：弗劳德数意义}
惯性力与重力之比。自由表面、水波、船模相似常用$Fr$。
\sect{问：欧拉方程}
无粘流体动量方程：$\rho D\bm V/Dt=-\nabla p+\rho\bm f$。Bernoulli可由其积分。
\sect{问：N--S方程}
不可压牛顿流体：$\rho D\bm V/Dt=-\nabla p+\mu\nabla^2\bm V+\rho\bm f$。粘性项导致扩散和能量损失。
\sect{问：H--P适用条件}
定常、不可压、牛顿流体、圆管、层流、充分发展。出口入口发展段、湍流不适用。
\sect{问：边界层基本思想}
高Re时粘性集中在壁面附近薄层，外部可按无粘流求解，边界层内满足无滑移。
\sect{问：分离原因}
逆压梯度使近壁流体动能不足而回流，壁面剪应力降为零后变号。
\sect{问：阻塞定义}
喷管最小截面达到$M=1$，质量流量最大，下游背压继续降低不再影响上游质量流量。
\sect{问：激波定义}
超声速压缩扰动形成的薄间断层，跨波静压、密度、温度突升，马赫数下降，总压损失。
\sect{问：膨胀波定义}
超声速流绕凸角膨胀形成的连续波扇，流动加速、压力温度下降，过程近似等熵。
\sect{问：Kutta条件}
实际有粘流体绕翼型后缘速度有限，确定环量大小，从而确定升力。
\sect{问：库塔--儒可夫斯基}
二维翼型单位展长升力$L'=\rho U\Gamma$，方向由来流和环量决定。
\sect{问：达朗贝尔佯谬}
理想无粘不可压定常绕流物体阻力为零，与现实不符，原因是忽略粘性分离。

\chap{M19 数值替换模板}
\sect{模板：先写符号后代数}
不要边想边代。先写$Re=Vd/\nu=\cdots$，再写数值。公式正确即使数值错也能保分。
\sect{模板：面积换算}
$A=\pi d^2/4=\pi(0.1)^2/4=7.85\times10^{-3}{\rm m^2}$。这类中间值可写在草稿区。
\sect{模板：速度换算}
$V=Q/A$。若$Q$给$m^3/h$，先除3600。若给L/s，乘$10^{-3}$。
\sect{模板：压强水头}
$p/(\rho g)=70000/(1000\times9.81)=7.14$m。kPa转水头直接除9.81。
\sect{模板：雷诺数}
$Re=Vd/\nu$。把$\nu$写成科学计数法，防止$10^{-6}$漏掉。
\sect{模板：损失水头}
$h_f=\lambda(L/d)V^2/(2g)$。先算$L/d$和速度头，再乘$\lambda$。
\sect{模板：功率}
$P=\rho gQH$，水中可速算$P({\rm kW})\approx9.81QH$，$Q$用$m^3/s$。
\sect{模板：等熵压比}
$p=p_0/(1+0.2M^2)^{3.5}$。如果$M=1$，直接$p=0.528p_0$。
\sect{模板：质量流量}
$\dot m=\rho VA$。可压缩出口：先$p,T$求$\rho=p/(RT)$，再$V=M\sqrt{\gamma RT}$。
\sect{模板：压力中心}
矩形竖直上边距液面$h_0$：$y_C=h_0+h/2$，$I_C=bh^3/12$，$y_D=y_C+I_C/(y_CA)$。
\sect{模板：Stokes}
终速$U=(\rho_s-\rho)gd^2/(18\mu)$。代完后算$Re=\rho Ud/\mu$检查。
\chap{M20 备用补充：冷门但可能考}
\sect{毛细现象}
毛细上升$h=2\sigma\cos\theta/(\rho gr)$。若玻璃管很细、题给表面张力和接触角，用此式。
\sect{表面张力压差}
球形液滴内外压差$\Delta p=2\sigma/R$；肥皂泡两界面$\Delta p=4\sigma/R$。
\sect{水击基本式}
快速关阀压升$\Delta p=\rho a\Delta V$。若课程提到水击，声速$a$和速度突变决定压强突变。
\sect{不可压能量损失}
损失代表机械能转化为内能，不能用伯努利恢复。局部压力可能恢复，但总水头下降。
\sect{涡量输运概念}
无粘无体力正压流中涡量随流体运动保持；粘性会扩散涡量。边界层是壁面产生涡量的重要区域。
\sect{Kelvin环量定理}
理想无粘、正压、保守体力下，随流体运动闭合曲线的环量守恒。解释势流中环量需由边界/初始条件决定。
\sect{控制体能量单位}
压力形式Pa：$p+\frac12\rho V^2+\rho gz$。水头形式m：$p/\rho g+V^2/2g+z$。两种不要混写。
\sect{开口容器边界}
自由液面$p=p_a$，若面积很大$V\approx0$。若容器运动或液面下降，$V$可能不可忽略。
\sect{流线弯曲压强}
曲线流动法向Euler：$\partial p/\partial n=\rho V^2/R$，压强向曲率中心外侧增大/按法向定义判断。可解释弯曲流线压力分布。
\sect{自由射流}
射流暴露于大气，边界静压等于大气压；射流内部若截面变化，速度由连续和能量决定。
\sect{扩压器效率}
若题给扩压器效率，实际压力升高=效率$\times$理想动压回收。无效率时按给定损失系数。
\sect{喷嘴效率}
实际出口动能=效率$\times$理想可用压降，或损失扣除后再求$V$。按题定义写。
\sect{多相/密度差}
U管、浮力、沉降都靠密度差。看清是$\rho_m-\rho$还是$\rho_s-\rho$。
\sect{非定常伯努利}
势流非定常保留$\partial\phi/\partial t$；U管振荡、液柱运动可能用到。定常题该项为0。
\sect{旋转流Bernoulli}
有旋流不能全场同一Bernoulli常数，只能沿流线用。自由涡外部虽局部无旋但有奇点，跨越奇点需小心。
\sect{二维流量单位}
二维势流的$\psi$常表示单位宽流量，单位$m^2/s$。总流量需乘宽度。
\sect{坐标选取}
管路题高度$z$向上为正；静水压深度$h$向下为正。两套符号不要混。
\sect{负号处理}
若压降写$\Delta p=p_1-p_2$，通常为正；若写$dp/dx$，沿流向压力下降时$dp/dx<0$。H--P常写$-dp/dx>0$。
\sect{图题入口：闸门/曲面}
看见闸门、铰链、拉索：第一行写$F=\rho gh_CA$；看见曲面：第一行写$F_x$投影面、$F_y=\rho gV_p$。后面一律对铰点取矩。
\sect{图题入口：大水箱+小孔}
大水面到孔口伯努利：$p_a+\rho gz_1=p_a+\rho gz_2+\frac12\rho V^2$，故$V=C_d\sqrt{2gH}$；变水位再接$A_t dH/dt=-A_oV$。
\sect{图题入口：弯管/喷嘴}
先连续$Q=AV$；再伯努利求速度或压强；最后动量$\sum\bm F=\rho Q(\bm V_2-\bm V_1)$求支座/管壁力，最后取反。
\sect{图题入口：管路+泵/水轮机}
选两个截面，写$p/\rho g+V^2/2g+z+H_p=p/\rho g+V^2/2g+z+H_T+h_f+\sum h_\zeta$。再由$Q,d$算$V,Re,\lambda$。
\sect{图题入口：收缩/扩张喷管}
题图有喉部：先判是否$M_t=1$。阻塞后$\dot m$由$p_0,T_0,A^*$定；出口压强再决定正激波、过膨胀或欠膨胀。
\sect{图题入口：平板边界层}
先算$Re_x=Ux/\nu$或$Re_L=UL/\nu$。层流用$\delta=5x/\sqrt{Re_x}$，全板阻力用$\bar C_f$乘动压和面积。
\sect{图题入口：圆柱/半圆柱势流}
先写$V_\theta=2U\sin\theta$，再$p=p_\infty+\frac12\rho(U^2-V_\theta^2)$。问升力/浮起就对压力分量积分或列竖向平衡。
\sect{图题入口：激波/膨胀角}
壁面向流体内折：斜激波；壁面向外展开：PM膨胀。斜激波用$M_{n1}=M_1\sin\beta$；膨胀用$\nu(M_2)-\nu(M_1)=\delta$。
\sect{模板：边界层阻力}
$Re_L=UL/\nu$，$\bar C_f$，$A$，$D=\bar C_f0.5\rho U^2A$。按这四步写。

\chap{M21 最后30秒题图决策卡}
\sect{图里有“自由液面”}
自由液面通常$p=p_a$；若水箱很大取$V\approx0$。若液面随时间动，必须接连续方程$A\,dh/dt=\pm Q$，不是只列伯努利。
\sect{图里有“压力表”}
表压进入能量方程可直接当$p/\rho g$；若题问汽蚀/蒸汽压，必须换绝压：$p_{\rm abs}=p_g+p_a$。
\sect{图里有“U形压差计”}
先写两连接点压力差。若测同一流体管路，常用$\Delta p=(\rho_m-\rho)g\Delta h$；若一端接大气，用$p-p_a=\rho_m g\Delta h$。
\sect{图里有“铰链+拉杆”}
别先求反力。先求水压力合力和压力中心，再对铰链取矩求拉力；最后若题要铰链反力，再用力平衡。
\sect{图里有“弯头/喷嘴出口”}
出口若射入大气，出口静压为大气压。支座力方向容易反：先求外界对流体力，再取相反即流体对管壁力。
\sect{图里有“泵在吸水侧”}
最高安装高度由泵入口最低允许绝压控制，不由泵功率直接控制。能量方程从水面列到泵入口。
\sect{图里有“水轮机”}
水轮机不是加能，是取能。能量式中$H_T$放在下游侧或从上游总水头中减去；功率$P=\eta\rho gQH_T$。
\sect{图里有“突然扩大/缩小”}
突扩损失常用$h=(V_1-V_2)^2/(2g)$；突缩或阀门按题给$\zeta V^2/(2g)$。速度取对应局部元件处速度。
\sect{图里有“圆管层流”}
H--P必须满足充分发展层流。若题给入口段/发展段，不要直接套抛物线；若给壁面剪切力，可由$\tau_w=d\Delta p/(4L)$反推压降。
\sect{图里有“平板零攻角”}
摩擦阻力用湿表面积，单面$Lb$，双面$2Lb$。空气和水同速同板，阻力差主要来自$\rho,\nu$导致的$Re$和$C_f$。
\sect{图里有“圆柱/球绕流”}
若是理想势流，压力积分阻力为0；若题给$C_D$，阻力$D=C_D(0.5\rho U^2)A_{\rm front}$，面积取迎风面积。
\sect{图里有“喉部”}
喉部是最小截面。只有阻塞时喉部$M=1$；没阻塞时喉部也可能$M<1$。不要见喉部就直接用0.528。
\sect{图里有“背压$p_b$”}
$p_b$是外界/下游压力，不一定等于喷管出口内压$p_e$。只有全亚声速收缩管或设计匹配等特殊状态才常有$p_e=p_b$。
\sect{图里有“正激波”}
波前必须超声速，波后亚声速。过波$p,T,\rho$升，$M$降，$T_0$不变，$p_0$下降；后面若继续等熵，用新的$p_{02}$。
\sect{图里有“斜面压缩角”}
超声速流遇凹角/压缩角，用斜激波；求解顺序是$M_1,\delta\to\beta\to M_{n1}\to$正激波关系$\to M_2$。
\sect{图里有“凸角膨胀”}
超声速流遇凸角，用PM膨胀扇。$\nu(M_2)=\nu(M_1)+\delta$，压力温度下降，总压不变。

\chap{M22 题图常见误判}
\sect{喷管出口线}
出口截面内压$p_e$不一定等于背压$p_b$；只有匹配或全亚声速等特殊情况才等。
\sect{喉部面积}
$A_t$是几何最小面积；$A^*$是临界面积。阻塞时二者相等，未阻塞时不能乱等。
\sect{平均速度}
管路能量题用截面平均速度$V=Q/A$；边界层/壁律题里的$u(y)$是局部速度。
\sect{表压零点}
开口水面表压为0，不是绝压为0；涉及汽化压、气体状态方程时必须用绝压。
\sect{方向符号}
静水深度$h$常向下为正，能量方程高度$z$向上为正。换题型时不要把符号带错。
\sect{摩擦面积}
平板阻力题问单面/双面要看图；管内沿程损失不用外表面积，用$L/d$和速度头。
\sect{局部损失速度}
不同管径时，每个阀门/弯头的$\zeta V^2/2g$用该元件所在管段速度。
\sect{激波查表}
正激波表输入$M_1$；斜激波先转成法向$M_{n1}$；膨胀表输入PM函数$\nu(M)$。
\sect{应力题}
求某平面应力矢量直接$P\bm n$；法向分量是投影$(P\bm n)\cdot\bm n$，不是矩阵对角元。
\end{multicols}
\end{document}
